%%%%%%%%%%%%%%%%%%%%%%%%%%%%%%%%%%%%%%%%%
% The Legrand Orange Book
% LaTeX Template
% Version 1.4 (12/4/14)
%
% This template has been downloaded from:
% http://www.LaTeXTembplates.com
%
% Original author:
% Mathias Legrand (legrand.mathias@gmail.com)
%
% License:
% CC BY-NC-SA 3.0 (http://creativecommons.org/licenses/by-nc-sa/3.0/)
%
% Compiling this template:
% This template uses biber for its bibliography and makeindex for its index.
% When you first open the template, compile it from the command line with the 
% commands below to make sure your LaTeX distribution is configured correctly:
%
% 1) pdflatex main
% 2) makeindex main.idx -s StyleInd.ist
% 3) biber main
% 4) pdflatex main x 2
%
% After this, when you wish to update the bibliography/index use the appropriate
% command above and make sure to compile with pdflatex several times 
% afterwards to propagate your changes to the document.
%
% This template also uses a number of packages which may need to be
% updated to the newest versions for the template to compile. It is strongly
% recommended you update your LaTeX distribution if you have any
% compilation errors.
%
% Important note:
% Chapter heading images should have a 2:1 width:height ratio,
% e.g. 920px width and 460px height.
%
%%%%%%%%%%%%%%%%%%%%%%%%%%%%%%%%%%%%%%%%%

%----------------------------------------------------------------------------------------
%	PACKAGES AND OTHER DOCUMENT CONFIGURATIONS
%----------------------------------------------------------------------------------------

\documentclass[11pt,fleqn]{book} % Default font size and left-justified equations

\usepackage[top=3cm,bottom=3cm,left=3.2cm,right=3.2cm,headsep=10pt,a4paper]{geometry} % Page margins

\usepackage{xcolor} % Required for specifying colors by name
\definecolor{ocre}{RGB}{243,102,25} % Define the orange color used for highlighting throughout the book
\usepackage{amsmath} %Required for some mathmatical fomulas 
\usepackage{enumitem} %lists
% Font Settings
\usepackage{avant} % Use the Avantgarde font for headings
%\usepackage{times} % Use the Times font for headings
\usepackage{mathptmx} % Use the Adobe Times Roman as the default text font together with math symbols from the Sym­bol, Chancery and Com­puter Modern fonts

\usepackage{bm}
\usepackage{bbm}
%\usepackage{parskip}

\usepackage{microtype} % Slightly tweak font spacing for aesthetics
\usepackage[utf8]{inputenc} % Required for including letters with accents
\usepackage[T1]{fontenc} % Use 8-bit encoding that has 256 glyphs
\usepackage{multicol}

% Bibliography
\usepackage[style=alphabetic,sorting=nyt,sortcites=true,autopunct=true,babel=hyphen,hyperref=true,abbreviate=false,backref=true,backend=biber]{biblatex}
\addbibresource{bibliography.bib} % BibTeX bibliography file
\defbibheading{bibempty}{}

% Index
\usepackage{calc} % For simpler calculation - used for spacing the index letter headings correctly
\usepackage{makeidx} % Required to make an index
\makeindex % Tells LaTeX to create the files required for indexing

%----------------------------------------------------------------------------------------

%fdlkjadlkjdfsal;kjafds;lkjafsd;lkjfsad

%%%%%%%%%%%%%%%%%%%%%%%%%%%%%%%%%%%%%%%%%
% The Legrand Orange Book
% Structural Definitions File
% Version 2.0 (9/2/15)
%
% Original author:
% Mathias Legrand (legrand.mathias@gmail.com) with modifications by:
% Vel (vel@latextemplates.com)
% 
% This file has been downloaded from:
% http://www.LaTeXTemplates.com
%
% License:
% CC BY-NC-SA 3.0 (http://creativecommons.org/licenses/by-nc-sa/3.0/)
%
%%%%%%%%%%%%%%%%%%%%%%%%%%%%%%%%%%%%%%%%%

%----------------------------------------------------------------------------------------
%	VARIOUS REQUIRED PACKAGES AND CONFIGURATIONS
%----------------------------------------------------------------------------------------

\usepackage[top=3cm,bottom=3cm,left=3cm,right=3cm,headsep=10pt,a4paper]{geometry} % Page margins

\usepackage{graphicx} % Required for including pictures
\graphicspath{{Pictures/}} % Specifies the directory where pictures are stored

\usepackage{lipsum} % Inserts dummy text

\usepackage{tikz} % Required for drawing custom shapes

\usepackage[english]{babel} % English language/hyphenation

\usepackage{enumitem} % Customize lists
\setlist{nolistsep} % Reduce spacing between bullet points and numbered lists

\usepackage{booktabs} % Required for nicer horizontal rules in tables

\usepackage{xcolor} % Required for specifying colors by name
\definecolor{ocre}{RGB}{243,102,25} % Define the orange color used for highlighting throughout the book

%----------------------------------------------------------------------------------------
%	FONTS
%----------------------------------------------------------------------------------------

\usepackage{avant} % Use the Avantgarde font for headings
%\usepackage{times} % Use the Times font for headings
\usepackage{mathptmx} % Use the Adobe Times Roman as the default text font together with math symbols from the Sym­bol, Chancery and Com­puter Modern fonts

\usepackage{microtype} % Slightly tweak font spacing for aesthetics
\usepackage[utf8]{inputenc} % Required for including letters with accents
\usepackage[T1]{fontenc} % Use 8-bit encoding that has 256 glyphs

%----------------------------------------------------------------------------------------
%	BIBLIOGRAPHY AND INDEX
%----------------------------------------------------------------------------------------

\usepackage[style=alphabetic,citestyle=numeric,sorting=nyt,sortcites=true,autopunct=true,babel=hyphen,hyperref=true,abbreviate=false,backref=true,backend=biber]{biblatex}
\addbibresource{bibliography.bib} % BibTeX bibliography file
\defbibheading{bibempty}{}

\usepackage{calc} % For simpler calculation - used for spacing the index letter headings correctly
\usepackage{makeidx} % Required to make an index
\makeindex % Tells LaTeX to create the files required for indexing

%----------------------------------------------------------------------------------------
%	MAIN TABLE OF CONTENTS
%----------------------------------------------------------------------------------------

\usepackage{titletoc} % Required for manipulating the table of contents

\contentsmargin{0cm} % Removes the default margin

% Part text styling
\titlecontents{part}[0cm]
{\addvspace{20pt}\centering\large\bfseries}
{}
{}
{}

% Chapter text styling
\titlecontents{chapter}[1.25cm] % Indentation
{\addvspace{12pt}\large\sffamily\bfseries} % Spacing and font options for chapters
{\color{ocre!60}\contentslabel[\Large\thecontentslabel]{1.25cm}\color{ocre}} % Chapter number
{\color{ocre}}  
{\color{ocre!60}\normalsize\;\titlerule*[.5pc]{.}\;\thecontentspage} % Page number

% Section text styling
\titlecontents{section}[1.25cm] % Indentation
{\addvspace{3pt}\sffamily\bfseries} % Spacing and font options for sections
{\contentslabel[\thecontentslabel]{1.25cm}} % Section number
{}
{\hfill\color{black}\thecontentspage} % Page number
[]

% Subsection text styling
\titlecontents{subsection}[1.25cm] % Indentation
{\addvspace{1pt}\sffamily\small} % Spacing and font options for subsections
{\contentslabel[\thecontentslabel]{1.25cm}} % Subsection number
{}
{\ \titlerule*[.5pc]{.}\;\thecontentspage} % Page number
[]

% List of figures
\titlecontents{figure}[0em]
{\addvspace{-5pt}\sffamily}
{\thecontentslabel\hspace*{1em}}
{}
{\ \titlerule*[.5pc]{.}\;\thecontentspage}
[]

% List of tables
\titlecontents{table}[0em]
{\addvspace{-5pt}\sffamily}
{\thecontentslabel\hspace*{1em}}
{}
{\ \titlerule*[.5pc]{.}\;\thecontentspage}
[]

%----------------------------------------------------------------------------------------
%	MINI TABLE OF CONTENTS IN PART HEADS
%----------------------------------------------------------------------------------------

% Chapter text styling
\titlecontents{lchapter}[0em] % Indenting
{\addvspace{15pt}\large\sffamily\bfseries} % Spacing and font options for chapters
{\color{ocre}\contentslabel[\Large\thecontentslabel]{1.25cm}\color{ocre}} % Chapter number
{}  
{\color{ocre}\normalsize\sffamily\bfseries\;\titlerule*[.5pc]{.}\;\thecontentspage} % Page number

% Section text styling
\titlecontents{lsection}[0em] % Indenting
{\sffamily\small} % Spacing and font options for sections
{\contentslabel[\thecontentslabel]{1.25cm}} % Section number
{}
{}

% Subsection text styling
\titlecontents{lsubsection}[.5em] % Indentation
{\normalfont\footnotesize\sffamily} % Font settings
{}
{}
{}

%----------------------------------------------------------------------------------------
%	PAGE HEADERS
%----------------------------------------------------------------------------------------

\usepackage{fancyhdr} % Required for header and footer configuration

\pagestyle{fancy}
\renewcommand{\chaptermark}[1]{\markboth{\sffamily\normalsize\bfseries\chaptername\ \thechapter.\ #1}{}} % Chapter text font settings
\renewcommand{\sectionmark}[1]{\markright{\sffamily\normalsize\thesection\hspace{5pt}#1}{}} % Section text font settings
\fancyhf{} \fancyhead[LE,RO]{\sffamily\normalsize\thepage} % Font setting for the page number in the header
\fancyhead[LO]{\rightmark} % Print the nearest section name on the left side of odd pages
\fancyhead[RE]{\leftmark} % Print the current chapter name on the right side of even pages
\renewcommand{\headrulewidth}{0.5pt} % Width of the rule under the header
\addtolength{\headheight}{2.5pt} % Increase the spacing around the header slightly
\renewcommand{\footrulewidth}{0pt} % Removes the rule in the footer
\fancypagestyle{plain}{\fancyhead{}\renewcommand{\headrulewidth}{0pt}} % Style for when a plain pagestyle is specified

% Removes the header from odd empty pages at the end of chapters
\makeatletter
\renewcommand{\cleardoublepage}{
\clearpage\ifodd\c@page\else
\hbox{}
\vspace*{\fill}
\thispagestyle{empty}
\newpage
\fi}

%----------------------------------------------------------------------------------------
%	THEOREM STYLES
%----------------------------------------------------------------------------------------

\usepackage{amsmath,amsfonts,amssymb,amsthm} % For math equations, theorems, symbols, etc

\newcommand{\intoo}[2]{\mathopen{]}#1\,;#2\mathclose{[}}
\newcommand{\ud}{\mathop{\mathrm{{}d}}\mathopen{}}
\newcommand{\intff}[2]{\mathopen{[}#1\,;#2\mathclose{]}}
\newtheorem{notation}{Notation}[chapter]

% Boxed/framed environments
\newtheoremstyle{ocrenumbox}% % Theorem style name
{0pt}% Space above
{0pt}% Space below
{\normalfont}% % Body font
{}% Indent amount
{\small\bf\sffamily\color{ocre}}% % Theorem head font
{\;}% Punctuation after theorem head
{0.25em}% Space after theorem head
{\small\sffamily\color{ocre}\thmname{#1}\nobreakspace\thmnumber{\@ifnotempty{#1}{}\@upn{#2}}% Theorem text (e.g. Theorem 2.1)
\thmnote{\nobreakspace\the\thm@notefont\sffamily\bfseries\color{black}---\nobreakspace#3.}} % Optional theorem note
\renewcommand{\qedsymbol}{$\blacksquare$}% Optional qed square

\newtheoremstyle{blacknumex}% Theorem style name
{5pt}% Space above
{5pt}% Space below
{\normalfont}% Body font
{} % Indent amount
{\small\bf\sffamily}% Theorem head font
{\;}% Punctuation after theorem head
{0.25em}% Space after theorem head
{\small\sffamily{\tiny\ensuremath{\blacksquare}}\nobreakspace\thmname{#1}\nobreakspace\thmnumber{\@ifnotempty{#1}{}\@upn{#2}}% Theorem text (e.g. Theorem 2.1)
\thmnote{\nobreakspace\the\thm@notefont\sffamily\bfseries---\nobreakspace#3.}}% Optional theorem note

\newtheoremstyle{blacknumbox} % Theorem style name
{0pt}% Space above
{0pt}% Space below
{\normalfont}% Body font
{}% Indent amount
{\small\bf\sffamily}% Theorem head font
{\;}% Punctuation after theorem head
{0.25em}% Space after theorem head
{\small\sffamily\thmname{#1}\nobreakspace\thmnumber{\@ifnotempty{#1}{}\@upn{#2}}% Theorem text (e.g. Theorem 2.1)
\thmnote{\nobreakspace\the\thm@notefont\sffamily\bfseries---\nobreakspace#3.}}% Optional theorem note

% Non-boxed/non-framed environments
\newtheoremstyle{ocrenum}% % Theorem style name
{5pt}% Space above
{5pt}% Space below
{\normalfont}% % Body font
{}% Indent amount
{\small\bf\sffamily\color{ocre}}% % Theorem head font
{\;}% Punctuation after theorem head
{0.25em}% Space after theorem head
{\small\sffamily\color{ocre}\thmname{#1}\nobreakspace\thmnumber{\@ifnotempty{#1}{}\@upn{#2}}% Theorem text (e.g. Theorem 2.1)
\thmnote{\nobreakspace\the\thm@notefont\sffamily\bfseries\color{black}---\nobreakspace#3.}} % Optional theorem note
\renewcommand{\qedsymbol}{$\blacksquare$}% Optional qed square
\makeatother

% Defines the theorem text style for each type of theorem to one of the three styles above
\newcounter{dummy} 
\numberwithin{dummy}{section}
\theoremstyle{ocrenumbox}
\newtheorem{theoremeT}[dummy]{Theorem}
\newtheorem{problem}{Problem}[chapter]
\newtheorem{exerciseT}{Exercise}[chapter]
\theoremstyle{blacknumex}
\newtheorem{exampleT}{Example}[chapter]
\theoremstyle{blacknumbox}
\newtheorem{vocabulary}{Vocabulary}[chapter]
\newtheorem{definitionT}{Definition}[section]
\newtheorem{corollaryT}[dummy]{Corollary}
\theoremstyle{ocrenum}
\newtheorem{proposition}[dummy]{Proposition}

%----------------------------------------------------------------------------------------
%	DEFINITION OF COLORED BOXES
%----------------------------------------------------------------------------------------

\RequirePackage[framemethod=default]{mdframed} % Required for creating the theorem, definition, exercise and corollary boxes

% Theorem box
\newmdenv[skipabove=7pt,
skipbelow=7pt,
backgroundcolor=black!5,
linecolor=ocre,
innerleftmargin=5pt,
innerrightmargin=5pt,
innertopmargin=5pt,
leftmargin=0cm,
rightmargin=0cm,
innerbottommargin=5pt]{tBox}

% Exercise box	  
\newmdenv[skipabove=7pt,
skipbelow=7pt,
rightline=false,
leftline=true,
topline=false,
bottomline=false,
backgroundcolor=ocre!10,
linecolor=ocre,
innerleftmargin=5pt,
innerrightmargin=5pt,
innertopmargin=5pt,
innerbottommargin=5pt,
leftmargin=0cm,
rightmargin=0cm,
linewidth=4pt]{eBox}	

% Definition box
\newmdenv[skipabove=7pt,
skipbelow=7pt,
rightline=false,
leftline=true,
topline=false,
bottomline=false,
linecolor=ocre,
innerleftmargin=5pt,
innerrightmargin=5pt,
innertopmargin=0pt,
leftmargin=0cm,
rightmargin=0cm,
linewidth=4pt,
innerbottommargin=0pt]{dBox}	

% Corollary box
\newmdenv[skipabove=7pt,
skipbelow=7pt,
rightline=false,
leftline=true,
topline=false,
bottomline=false,
linecolor=gray,
backgroundcolor=black!5,
innerleftmargin=5pt,
innerrightmargin=5pt,
innertopmargin=5pt,
leftmargin=0cm,
rightmargin=0cm,
linewidth=4pt,
innerbottommargin=5pt]{cBox}

% Creates an environment for each type of theorem and assigns it a theorem text style from the "Theorem Styles" section above and a colored box from above
\newenvironment{theorem}{\begin{tBox}\begin{theoremeT}}{\end{theoremeT}\end{tBox}}
\newenvironment{exercise}{\begin{eBox}\begin{exerciseT}}{\hfill{\color{ocre}\tiny\ensuremath{\blacksquare}}\end{exerciseT}\end{eBox}}				  
\newenvironment{definition}{\begin{dBox}\begin{definitionT}}{\end{definitionT}\end{dBox}}	
\newenvironment{example}{\begin{exampleT}}{\hfill{\tiny\ensuremath{\blacksquare}}\end{exampleT}}		
\newenvironment{corollary}{\begin{cBox}\begin{corollaryT}}{\end{corollaryT}\end{cBox}}	

%----------------------------------------------------------------------------------------
%	REMARK ENVIRONMENT
%----------------------------------------------------------------------------------------

\newenvironment{remark}{\par\vspace{10pt}\small % Vertical white space above the remark and smaller font size
\begin{list}{}{
\leftmargin=35pt % Indentation on the left
\rightmargin=25pt}\item\ignorespaces % Indentation on the right
\makebox[-2.5pt]{\begin{tikzpicture}[overlay]
\node[draw=ocre!60,line width=1pt,circle,fill=ocre!25,font=\sffamily\bfseries,inner sep=2pt,outer sep=0pt] at (-15pt,0pt){\textcolor{ocre}{R}};\end{tikzpicture}} % Orange R in a circle
\advance\baselineskip -1pt}{\end{list}\vskip5pt} % Tighter line spacing and white space after remark

%----------------------------------------------------------------------------------------
%	SECTION NUMBERING IN THE MARGIN
%----------------------------------------------------------------------------------------

\makeatletter
\renewcommand{\@seccntformat}[1]{\llap{\textcolor{ocre}{\csname the#1\endcsname}\hspace{1em}}}                    
\renewcommand{\section}{\@startsection{section}{1}{\z@}
{-4ex \@plus -1ex \@minus -.4ex}
{1ex \@plus.2ex }
{\normalfont\large\sffamily\bfseries}}
\renewcommand{\subsection}{\@startsection {subsection}{2}{\z@}
{-3ex \@plus -0.1ex \@minus -.4ex}
{0.5ex \@plus.2ex }
{\normalfont\sffamily\bfseries}}
\renewcommand{\subsubsection}{\@startsection {subsubsection}{3}{\z@}
{-2ex \@plus -0.1ex \@minus -.2ex}
{.2ex \@plus.2ex }
{\normalfont\small\sffamily\bfseries}}                        
\renewcommand\paragraph{\@startsection{paragraph}{4}{\z@}
{-2ex \@plus-.2ex \@minus .2ex}
{.1ex}
{\normalfont\small\sffamily\bfseries}}

%----------------------------------------------------------------------------------------
%	PART HEADINGS
%----------------------------------------------------------------------------------------

% numbered part in the table of contents
\newcommand{\@mypartnumtocformat}[2]{%
\setlength\fboxsep{0pt}%
\noindent\colorbox{ocre!20}{\strut\parbox[c][.7cm]{\ecart}{\color{ocre!70}\Large\sffamily\bfseries\centering#1}}\hskip\esp\colorbox{ocre!40}{\strut\parbox[c][.7cm]{\linewidth-\ecart-\esp}{\Large\sffamily\centering#2}}}%
%%%%%%%%%%%%%%%%%%%%%%%%%%%%%%%%%%
% unnumbered part in the table of contents
\newcommand{\@myparttocformat}[1]{%
\setlength\fboxsep{0pt}%
\noindent\colorbox{ocre!40}{\strut\parbox[c][.7cm]{\linewidth}{\Large\sffamily\centering#1}}}%
%%%%%%%%%%%%%%%%%%%%%%%%%%%%%%%%%%
\newlength\esp
\setlength\esp{4pt}
\newlength\ecart
\setlength\ecart{1.2cm-\esp}
\newcommand{\thepartimage}{}%
\newcommand{\partimage}[1]{\renewcommand{\thepartimage}{#1}}%
\def\@part[#1]#2{%
\ifnum \c@secnumdepth >-2\relax%
\refstepcounter{part}%
\addcontentsline{toc}{part}{\texorpdfstring{\protect\@mypartnumtocformat{\thepart}{#1}}{\partname~\thepart\ ---\ #1}}
\else%
\addcontentsline{toc}{part}{\texorpdfstring{\protect\@myparttocformat{#1}}{#1}}%
\fi%
\startcontents%
\markboth{}{}%
{\thispagestyle{empty}%
\begin{tikzpicture}[remember picture,overlay]%
\node at (current page.north west){\begin{tikzpicture}[remember picture,overlay]%	
\fill[ocre!20](0cm,0cm) rectangle (\paperwidth,-\paperheight);
\node[anchor=north] at (4cm,-3.25cm){\color{ocre!40}\fontsize{220}{100}\sffamily\bfseries\@Roman\c@part}; 
\node[anchor=south east] at (\paperwidth-1cm,-\paperheight+1cm){\parbox[t][][t]{8.5cm}{
\printcontents{l}{0}{\setcounter{tocdepth}{1}}%
}};
\node[anchor=north east] at (\paperwidth-1.5cm,-3.25cm){\parbox[t][][t]{15cm}{\strut\raggedleft\color{white}\fontsize{30}{30}\sffamily\bfseries#2}};
\end{tikzpicture}};
\end{tikzpicture}}%
\@endpart}
\def\@spart#1{%
\startcontents%
\phantomsection
{\thispagestyle{empty}%
\begin{tikzpicture}[remember picture,overlay]%
\node at (current page.north west){\begin{tikzpicture}[remember picture,overlay]%	
\fill[ocre!20](0cm,0cm) rectangle (\paperwidth,-\paperheight);
\node[anchor=north east] at (\paperwidth-1.5cm,-3.25cm){\parbox[t][][t]{15cm}{\strut\raggedleft\color{white}\fontsize{30}{30}\sffamily\bfseries#1}};
\end{tikzpicture}};
\end{tikzpicture}}
\addcontentsline{toc}{part}{\texorpdfstring{%
\setlength\fboxsep{0pt}%
\noindent\protect\colorbox{ocre!40}{\strut\protect\parbox[c][.7cm]{\linewidth}{\Large\sffamily\protect\centering #1\quad\mbox{}}}}{#1}}%
\@endpart}
\def\@endpart{\vfil\newpage
\if@twoside
\if@openright
\null
\thispagestyle{empty}%
\newpage
\fi
\fi
\if@tempswa
\twocolumn
\fi}

%----------------------------------------------------------------------------------------
%	CHAPTER HEADINGS
%----------------------------------------------------------------------------------------

\newcommand{\thechapterimage}{}%
\newcommand{\chapterimage}[1]{\renewcommand{\thechapterimage}{#1}}%
\def\@makechapterhead#1{%
{\parindent \z@ \raggedright \normalfont
\ifnum \c@secnumdepth >\m@ne
\if@mainmatter
\begin{tikzpicture}[remember picture,overlay]
\node at (current page.north west)
{\begin{tikzpicture}[remember picture,overlay]
\node[anchor=north west,inner sep=0pt] at (0,0) {\includegraphics[width=\paperwidth]{\thechapterimage}};
\draw[anchor=west] (\Gm@lmargin,-9cm) node [line width=2pt,rounded corners=15pt,draw=ocre,fill=white,fill opacity=0.5,inner sep=15pt]{\strut\makebox[22cm]{}};
\draw[anchor=west] (\Gm@lmargin+.3cm,-9cm) node {\huge\sffamily\bfseries\color{black}\thechapter. #1\strut};
\end{tikzpicture}};
\end{tikzpicture}
\else
\begin{tikzpicture}[remember picture,overlay]
\node at (current page.north west)
{\begin{tikzpicture}[remember picture,overlay]
\node[anchor=north west,inner sep=0pt] at (0,0) {\includegraphics[width=\paperwidth]{\thechapterimage}};
\draw[anchor=west] (\Gm@lmargin,-9cm) node [line width=2pt,rounded corners=15pt,draw=ocre,fill=white,fill opacity=0.5,inner sep=15pt]{\strut\makebox[22cm]{}};
\draw[anchor=west] (\Gm@lmargin+.3cm,-9cm) node {\huge\sffamily\bfseries\color{black}#1\strut};
\end{tikzpicture}};
\end{tikzpicture}
\fi\fi\par\vspace*{270\p@}}}

%-------------------------------------------

\def\@makeschapterhead#1{%
\begin{tikzpicture}[remember picture,overlay]
\node at (current page.north west)
{\begin{tikzpicture}[remember picture,overlay]
\node[anchor=north west,inner sep=0pt] at (0,0) {\includegraphics[width=\paperwidth]{\thechapterimage}};
\draw[anchor=west] (\Gm@lmargin,-9cm) node [line width=2pt,rounded corners=15pt,draw=ocre,fill=white,fill opacity=0.5,inner sep=15pt]{\strut\makebox[22cm]{}};
\draw[anchor=west] (\Gm@lmargin+.3cm,-9cm) node {\huge\sffamily\bfseries\color{black}#1\strut};
\end{tikzpicture}};
\end{tikzpicture}
\par\vspace*{270\p@}}
\makeatother

%----------------------------------------------------------------------------------------
%	HYPERLINKS IN THE DOCUMENTS
%----------------------------------------------------------------------------------------

\usepackage{hyperref}
\hypersetup{hidelinks,backref=true,pagebackref=true,hyperindex=true,colorlinks=false,breaklinks=true,urlcolor= ocre,bookmarks=true,bookmarksopen=false,pdftitle={Title},pdfauthor={Author}}
\usepackage{bookmark}
\bookmarksetup{
open,
numbered,
addtohook={%
\ifnum\bookmarkget{level}=0 % chapter
\bookmarksetup{bold}%
\fi
\ifnum\bookmarkget{level}=-1 % part
\bookmarksetup{color=ocre,bold}%
\fi
}
}
 % Insert the commands.tex file which contains the majority of the structure behind the template

\newlength\tindent
\setlength{\tindent}{\parindent}
\setlength{\parindent}{0pt}
\renewcommand{\indent}{\hspace*{\tindent}}

\begin{document}

%----------------------------------------------------------------------------------------
%	TITLE PAGE
%----------------------------------------------------------------------------------------

\begingroup
\thispagestyle{empty}
\AddToShipoutPicture*{\put(6,5){\includegraphics[scale=1]{background}}} % Image background
\centering
\vspace*{9cm}
\par\normalfont\fontsize{35}{35}\sffamily\selectfont
STAT 514 Lecture Notes\par % Book title
\vspace*{1cm}
{\Huge Dr. David Hunter}\par % Author name
\endgroup

%----------------------------------------------------------------------------------------
%	COPYRIGHT PAGE
%----------------------------------------------------------------------------------------

\newpage
~\vfill
\thispagestyle{empty}

\noindent Copyright \copyright\ 2014 John Smith\\ % Copyright notice

\noindent \textsc{Published by Publisher}\\ % Publisher

\noindent \textsc{book-website.com}\\ % URL

\noindent Licensed under the Creative Commons Attribution-NonCommercial 3.0 Unported License (the ``License''). You may not use this file except in compliance with the License. You may obtain a copy of the License at \url{http://creativecommons.org/licenses/by-nc/3.0}. Unless required by applicable law or agreed to in writing, software distributed under the License is distributed on an \textsc{``as is'' basis, without warranties or conditions of any kind}, either express or implied. See the License for the specific language governing permissions and limitations under the License.\\ % License information

\noindent \textit{First printing, March 2013} % Printing/edition date

%----------------------------------------------------------------------------------------
%	TABLE OF CONTENTS
%----------------------------------------------------------------------------------------

\chapterimage{chapter_head_1.pdf} % Table of contents heading image

\pagestyle{empty} % No headers

\tableofcontents % Print the table of contents itself

\cleardoublepage % Forces the first chapter to start on an odd page so it's on the right

\pagestyle{fancy} % Print headers again

%----------------------------------------------------------------------------------------
%	CHAPTER 1
%----------------------------------------------------------------------------------------

\chapterimage{chapter_head_2.pdf} % Chapter heading image

\chapter{Point Estimation - Chapter 7 Casella \& Berger}

\section{Introduction}\index{Introduction}

	In the simplest case, we have $n$ observations of data that we believe follow the same distribution.
	$$
		X_1, \dots, X_n \stackrel{iid}{\sim} f_\theta(x)
	$$

	where $f_\theta(x)$ is a density function involving a parameter $\theta$. Our goal is to learn something about $\theta$, which could be real or vector valued.\\

	\begin{definition}[Estimator]
		An \emph{estimator} of $\theta$ is any function $W(X_1, \dots, X_n)$ of the data. That is, an estimator is a \emph{statistic}.
	\end{definition}

	Note:
	\begin{enumerate}
		\item $W(\bm{X})$ may not depend on $\theta$.
		\item $W(\bm{X})$ should resemble or ``be close'' to $\theta$.
		\item An estimator is \emph{random}.
		\item $W(X_1, \dots, X_n)$ is the estimator, $W(x_1, \dots, x_n)$ is the fixed estimate.
	\end{enumerate}

	\begin{example}
		Suppose we have $n$ observations from an exponential distribution,
		$$
			X_1, \dots, X_n \stackrel{iid}{\sim} f_\theta(x)=\frac{1}{\theta}\exp\left\{-\frac{x}{\theta}\right\}\mathbbm{1}\{x>0\}
		$$
		for some $\theta > 0$. The \textbf{likelihood function} is equivalent to the joint density function, expressed as a function of $\theta$ rather than the data:
		$$
			L(\theta) = \prod_{i=1}^n \frac{1}{\theta}\exp\left\{-\frac{x}{\theta}\right\} = \frac{1}{\theta^n}\exp\left\{-\frac{1}{\theta}\sum_{i=1}^n x_i\right\}
		$$
		This function represents the \emph{likelihood} of observing the data we observed assuming the parameter was a particular value of $\theta$. If we can maximize this function, we can determine the $\hat{\theta}$ for which the likelihood of observing $\bm{X}$ was the highest. This might tell us something about the true value of $\theta$.

		\indent To maximize $L(\theta)$, we want to take the derivative, set it equal to 0, and solve for $\theta$. However, in many cases taking the derivative of the likelihood function will be very hard, if not impossible. We can use the fact that taking the logarithm does not change the location of extrema. The \textbf{log-likelihood function} in this case is
		\begin{equation*}
			\ell(\theta) = \log L(\theta) = -n \log \theta - \frac{1}{\theta} \sum_{i=1}^n x_i
		\end{equation*}
		Take the derivative with respect to the parameter and set equal to 0:
		\begin{eqnarray*}
			\ell'(\theta) &=& -\frac{n}{\theta} + \frac{1}{\theta^2} \sum_{i=1}^n x_i \enspace \stackrel{\textrm{set}}{=} \enspace 0\\
			\hat{\theta} &=& \frac{1}{n} \sum_{i=1}^n x_i
		\end{eqnarray*}
		Here $\hat{\theta}$ is an estimator (the sample mean). Since it maximizes $L(\theta)$, we call it the \textbf{maximum likelihood estimator} (MLE).
	\end{example}

\section{Methods of Finding Estimators}\index{Finding Estimators}

 \textbf{Setup:} $x_1, \dots, x_n \overset{iid}{\sim} f_\theta, \theta \in \Theta$  \\
 \textbf{Goal:} Estimate $\theta$, or some function $\tau(\theta)$\\

 \begin{definition}[Point Estimator]
 	A point estimator $T(x_1, \dots, x_n)$ is a function of the random sample. T should not depend on $\theta$. 
 \end{definition}

 \begin{remark}
 	Estimator is defined above, while an estimate is a fixed value. ($T(x_1 = x_1, \dots, x_n = x_n$)
 \end{remark}

 \subsection{Method of Moments}\index{Method of Moments}

	This method is one of the oldest, simplest estimators developed by Karl Pearson.\\ 
	\\
	\textbf{Construct:} $\begin{cases}
	               \frac{1}{n} \sum X_i = E_\theta(x) = \int_\theta x f_\theta(x) dx\\
	               \frac{1}{n} \sum X_i^2 = E_\theta(x^2) = \int_\theta x^2 f_\theta(x) dx\\
	               \vdots
	            \end{cases}$

	So, $\hat{\theta}_MM$ depends only on ($\frac{1}{n} \sum X_i = E_\theta(x), \dots, \frac{1}{n} \sum X_i = E_\theta(x^k)$).

	\begin{example}
		$x_1, \dots, x_n \overset{iid}{\sim} N(\mu, \sigma^2)$\\
		$\mu, \sigma^2$ unknown: $\Theta = \mathbb{R} x \mathbb{R}^+$\\
		\begin{align*}
		\frac{1}{n} \sum X_i = \mu &\implies \hat{\mu_{MM}} = \frac{1}{n} \sum X_i \triangleq \bar{X}\\
		\frac{1}{n} \sum X_i^2 = \sigma^2 + \mu^2 &\implies \hat{\sigma_MM^2} = \frac{1}{n} \sum (X_i - \bar{X})^2
		\end{align*}
	\end{example}

	\begin{example}
		$x_1, \dots, x_n \overset{iid}{\sim} \text{Bin}(k,p)$ where $k, p$ are unknown: $\Theta: \mathbb{N} \text{ x } [0,1]$.\\

		\begin{align*}
		\frac{1}{n} \sum X_i &= kp \triangleq \bar{X}\\
		\frac{1}{n} \sum X_i^2 &= kp(1-p) + k^2p^2 = \bar{X}(1-p) + \bar{X}^2
		\end{align*}
	\end{example}

	\begin{remark}
		Read about Satterthwait approximation
	\end{remark}

\subsection{Maximum Likelihood Estimation}\index{Maximum Likelihood Estimation}

	MLE makes data the most likely under the model specified. The quality of MLE depends on the quality of the model. Unlike Method of Moments, MLE does not lie outside of $\Theta$. \\
	\\
	Invariance Principle: if $\hat{\theta}$ is MLE for $\theta$, then $\tau(\hat{\theta})$ is MLE for $\tau(\theta)$. \\

	Roughly speaking, MLEs enjoy a nice property called \textbf{asymptotic efficiency}. They acheive the smallest possible varience in large sample. Similar to Lower Bound of Cramer-Rao as n -> infinity. 

	\begin{definition}[Score Function]
		Suppose that throughout $X_1,\dots,X_n \sim f_\theta(x)$ the \textbf{score function} is $\sum \frac{\delta}{\delta \theta} \log f_\theta(x_i)$. 
	\end{definition}

	Take a first-order Taylor expansion of $\ell(\theta)$ around true value ($\theta$). 
	\begin{align*}
	\ell^\prime(\hat{\theta})&\simeq \ell^\prime(\theta_0) + (\hat{\theta}-\theta_0)\ell^{\prime\prime}(\theta_0)\\
	\hat{\theta}-\theta_0&\simeq [\ell^{\prime\prime}(\theta_0)^-1][\ell^\prime(\hat{\theta})-\ell^\prime(\theta_0)]\\
		&=[\frac{-1}{n}\ell^{\prime\prime}(\theta_0)^{-1}][\frac{1}{n}\ell^\prime(\theta_0)]\\
	\frac{(\hat{\theta}-\theta_0)}{\sqrt{n}}&\simeq \frac{\frac{-1}{\sqrt{n}}\ell^\prime(\theta_0)}{\frac{-1}{n}\ell^{\prime\prime}(\theta_0)}\\
	\frac{1}{n}\ell^\prime(\theta_0)&= \frac{1}{n}\sum \frac{\delta}{\delta\theta}\log f_\theta(x)\\
	\frac{-1}{n}\ell^{\prime\prime}(\theta_0)&= \frac{1}{n}\sum \frac{\delta^2}{\delta\theta^2}\log f_\theta(x)|_{\theta=\theta_0}\\
	\end{align*}

	Thus, $\sqrt{n}(\hat{\theta}-\theta_0)\sim N(0,[I(\theta_0)]^{-1})$\\

	\textbf{Finding the Maxima:} $\theta \in \mathbb{R}^k$\\
	Suppose L is differentialbe, then the only possible candidtes for $\hat{\theta_{MLE}}$ are the stationary points of L which satisfy $\frac{\delta L}{\delta \theta_i} = 0, i=1, \dots, k$, which is a necessary condition for optimality.

	Suppose L is twice differentiable,

	\begin{center}
	H $\triangleq \begin{bmatrix}
	    \frac{\delta^2 L}{\delta \theta_1^2} & \frac{\delta^2 L}{\delta \theta_1 \delta \theta_2} & \dots & \frac{\delta^2 L}{\delta \theta_1 \delta \theta_k} \\
	    \vdots & \vdots & \ddots & \vdots \\
	        \frac{\delta^2 L}{\delta \theta_k \theta_1} & \dots & \dots & \frac{\delta^2 L}{\delta \theta_k^2}

	\end{bmatrix}$	
	\end{center}

	If $H|_{\theta = \hat{\theta_{ML}}}$ is negative definite, then $\hat{theta_{ML}}$ is the MLE. 

	\begin{definition}[Negative Definite]
		A matrix, $A_{nxn}$ is called \textbf{negative definite} if for any $(a_1, a_2, dots, a_n) \in \mathbb{R}^m$, $\sum\limits^n_{i,j = 1} a_i a_j A_{ij} < 0$. Equivalent characterization is the eigenvalues of A are negative.
	\end{definition}

	\begin{remark}
	More reading: C\&B sections 3.4, 7.2.3.
	\end{remark}

	\begin{remark}
		MLE and MM should match on: Bernoulli, Binomial, Normal, and Poisson distrabutions. 
	\end{remark}
%end subsection - MLE

\subsection{Bayes' Estimation}\index{Baye's Estimation}

	\begin{theorem}[Bayes' Theorem]
	$$\pi(\theta|x) = \frac{f(x|\theta)\pi(\theta)}{m(x)} $$
	LHS: posterior density of the parameters given data. This is what we want to know. The RHS includes
	\begin{enumerate}[label = \roman*)]
		\item likelihood
		\item prior density
		\item marginal density of x
	\end{enumerate}
		
	\end{theorem}

	\begin{proposition}[Dave's Bayes Way]
	 posterior $\propto$ prior*likelihood
	\end{proposition}

	\begin{remark}
		Bayesian estimators regularize MLE.
	\end{remark}

	\begin{definition}[Conjugate Family]
	If the prior is chosen from a family of distributions in such a way that the posterior is also a member, then we call this family of priors a \textbf{conjugate family}.
	\end{definition}


		\begin{table}[h]
		\centering
		\begin{tabular}{l l}
		\toprule
		\textbf{Likelihood} & \textbf{Conjugate Prior}\\
		\midrule
		Binom(p) & Beta($\alpha, \beta$) \\
		Poisson($\lambda$) & Gamma($\alpha, \beta$) \\
		Exp($\lambda$) & Gamma($\alpha, \beta$) \\
		Normal($\mu, \sigma^2$) & Gamma($\alpha, \beta$) \\
		\bottomrule
		\end{tabular}
		\end{table}

%end subsection - Baye's Estimator
%end section - methods of finding estimators

\section{Evaluating Estimation}\index{Evaluating Estimation}

\subsection{Mean Squared Error}\index{Mean Squared Error}

	\begin{remark} 
		Read Casella \& Berger Chapter 7.3 - Methods of Evaluating Estimation
	\end{remark} 


	\begin{definition}[Mean Squared Error]

	If $W(\bm{X})$ is an estimator of $\theta$, then the \textbf{mean squared error} (MSE) is defined as
	\[
		E_\theta\left[(W(\bm{X}) - \theta)^2\right].
	\]

	\end{definition}

	\begin{definition}[Unbiased estimator]

	If $W(\bm{X})$ is an estimator of $\theta$, we say that $W(\bm{X})$ is \textbf{unbiased} if
	\[
		E_\theta[W(\bm{X})] = \theta \quad \forall \theta.
	\]
	Furthermore, the \textbf{bias} of $W(\bm{X})$ is
	\[
		E_\theta[W(\bm{X})] - \theta.
	\]

	\end{definition}

	\begin{example}
		For $\theta$ > 0, let
		$$
			X_1, \dots, X_n \stackrel{iid}{\sim} f_\theta(x)=\theta x^{-2} \mathbbm{1}\{x>\theta\}
		$$
	Find the MLE of $\theta$.
	$$
		L(\theta)=\theta^n \prod_{i=1}^n x_i^{-2} \prod_{i=1}^n \mathbbm{1}\{x>\theta\}$$
		\\
		$$\hat{\theta}=\text{minimum of }x_i
	$$
	\end{example}

	\begin{theorem}
	$\textrm{MSE}(W) = \textrm{Bias}^2 + \textrm{Var}(W)$
	\end{theorem}

	\begin{proof}
	\begin{eqnarray*}
		E[(W(\bm{X})-\theta)^2] &=& E[(W-E[W]+E[W]-\theta)^2]\\
								&=& E[(W-E[W])^2]+E[(E[W]-\theta)^2]+2E[(W-E[W])(E[W]-\theta)]\\
								&=& \textrm{Var}(W) + \textrm{Bias}^2(W) + 0
	\end{eqnarray*}	
	\end{proof}
	
	\begin{example}
		$X_1, \dots, X_n \stackrel{iid}{\sim} N(\mu, \sigma^2)$ where $\sigma^2$ is known.
		$\hat{\mu}_{ML} = \bar{X}$; and $\hat{\mu} = X_1$

		\begin{multicols}{2}
		  \begin{align*}
		    \text{MSE}(\hat{\mu}_{ML}) &= [E(\hat{\mu}_{ML} - \mu)]^2 + \text{Var}(\bar{X})\\
		    	&= E(\bar{X}) - \mu + \frac{\sigma^2}{n}\\
		    	&= \frac{\sigma^2}{n}
		  \end{align*}\break
		  \begin{align*}
		    \text{MSE}(\hat{\mu}) &= [E(\hat{\mu} - \mu)]^2 + \text{Var}(X_1)\\
		    	&= E(X_1) - \mu + \sigma^2\\
		    	&= \sigma^2
		  \end{align*}
		\end{multicols}

		So, for n > 1, $\text{MSE}(\hat{\mu}_{ML})$ is better than $\text{MSE}(\hat{\mu})$.
	\end{example}

\subsection{Best Unbiased Estimator}\index{Best Unbiased Estimator}
		What does "best" mean? Answer: Minimum variance. \\
	
	While we can NOT find estimators that have smallest MSE for and $\theta$, we CAN find estimators within the class of unbiased estimators that achieve smallest MSE for all $\theta \in \Theta$.\\

	\begin{definition}[Best Unbiased Estimator]
		An estimator W* is a \textbf{best unbiased estimator}* of $\tau(\theta)$ if it satisties
		\[
			E_\theta (W*) = \tau(\theta)
		\] 
		for all $\theta$ and, for any other estimator W with 
		\[
			E_\theta (W) = \tau(\theta)
		\]
		we have 
		\[
			\text{Var}_\theta (W*) \le \text{Var}_\theta (W)
		\]
		for all $\theta$. W* is also called a $\textbf{uniform minimum variance unbiased estimator}$ (UMVUE) of $\tau(\theta)$.
	\end{definition}

	\textbf{Main Result:} Under some assumptions we can establish a lower bound on Var(W(X)). The question of how small this variance could be is answered by Cramer-Rao Inequality.

	\begin{remark}
		Spoilers: We will show that there are two main route for finding the UMVUE. 
		\begin{enumerate}
			\item Check if estimator attains the Cramer-Rao Lower Bound. 
			\item Find a complete and sufficient statistic then a function of it that is unbiased. 
			\begin{enumerate}
				\item Check to see if density if part of the exponential family. T is complete and sufficient.
				\item Use factorization theorem to find a sufficient statistic then the definition of completeness and Lehmann-Sheffe to get UMVUE
			\end{enumerate}
		\end{enumerate}
	\end{remark}

	But first, we need the following:

	\begin{theorem}[Cauchy-Scwarz Inequality]
		$$||<x,y>||^2 \le ||<x,x>|| * ||<y,y>||$$
	In terms of E(X) if A \& B have $\mu$ = 0:
	$$(E(AB))^2 \le E(A^2) * E(B^2)$$

	Or in terms of covariance:
	$$Cov^2(AB) \le Var(A)*Var(B)$$
	\end{theorem}
	
	\begin{proof}
	Let $D = B - \frac{E(AB)}{E(A^2)}A$, given $D^2 \ge 0$.
		\begin{align*}
			E(D^2) &= E(B^2 - 2(B)(\frac{E(AB)}{E(A^2)}A)+(\frac{E(AB)}{E(A^2)}A)^2)\\
				&= E(B^2)-2\left( \frac{E(AB)^2}{E(A^2)} \right) + \frac{E(AB)}{E(A^2)}E(A^2)\\
				&= E(B^2) - \frac{E(AB)^2}{E(A^2)} \ge 0
		\end{align*}	
	\end{proof}

	\begin{theorem}[Cramer-Rao Inequality]
		Also: Information Inequality. 
		$$\text{Var}(W(\underline{X})) \ge \frac{(\Psi^\prime(\theta))^2}{I(\theta)}$$
		where, 
		$\Psi(\theta)=E_\theta(W(\underline{X}))$\\
		and, 
		the Fisher/Expected information, $I(\theta)=E\left((\frac{\delta}{\delta \theta}\log f_\theta(\underline{X}))^2 \right)$.

		
	\end{theorem}

	\begin{proof}
		Follows from Cauchy-Schwarz Inequality 

		$$Cov(W(\underline{x}), \frac{\delta}{\delta\theta}\log f_\theta(\underline{X})))^2 \le \text{Var}(W(\underline{X}))*\text{Var}(\frac{\delta}{\delta\theta}\log f_\theta(\underline{X})) $$

	\end{proof}

	\begin{definition}[Regular]
		A family, $\mathcal{F} = \{f_\theta:\theta \in \Theta\}$ is called \textbf{regular} if:
		\begin{enumerate}[label = (\Roman*)]
			\item The set $A = \{x: f_\theta (x) > 0\}$ does not depend on $\theta$.  Forall $x\in A, \theta \in \Theta, \frac{\delta}{\delta\theta} \log(f_\theta(x))$ exists and is finite.
			\item For any estimator W such that $E_\theta|W| < \infty$ for all $\theta \in \Theta$, the operations of integration and differentiation by $\theta$ can be interchanged when the r.h.s is finite.
		\end{enumerate}
	\end{definition}

	\textbf{Assumptions:}
	\begin{enumerate}
		\item I($\theta)=\text{E}_\theta \left((\frac{\delta}{\delta \theta}\log f_\theta(\underline{X}))^2\right) = \text{Var}\left((\frac{\delta}{\delta \theta}\log f_\theta(\underline{X})^2)\right)$ is well defined and I($\theta)>0$.
		\item $E \left((\frac{\delta}{\delta \theta}\log f_\theta(\underline{X}))^2 \right)=0$. Thus, I($\theta$) is just varience.
		\item $E \left((W(\underline{X})\frac{\delta}{\delta \theta}\log f_\theta(\underline{X}))^2 \right)=\Psi^\prime(\theta)$
	\end{enumerate}

	So, \\
	$Cov(W(\underline{x}), \frac{\delta}{\delta\theta}\log f_\theta(\underline{X})))^2 \le \text{Var}(W(\underline{X}))*\text{Var}(\frac{\delta}{\delta\theta}\log f_\theta(\underline{X})) \simeq \text{Var}(W(\underline{X})) \ge \frac{(\Psi^\prime(\theta))^2}{I(\theta)}$

	\begin{remark}
		\begin{enumerate}[label = (\roman*)]
			\item Suppose W(x) is an unbiased estimator of $\theta$. Then $(\Psi^\prime(\theta))^2 = 1$ and the bound has nothing to do with the estimator anymore, but only deals with $f_\theta$.
			\item Suppose $X_1, \dots, X_n \stackrel{iid}{\sim} f_\theta$, then $I(\theta) = nI(\theta)$ where the latter information is related to only a single variable, X.
			\item This result also holds for discrete distributions (as long as the summation and differentiation are interchangeable).
			\item The CR bound might be stictly smaller than the varience of any estimator.
			\item Under additional assumptions on $f_\theta$, 
			$$I(\theta) = -E\left[\frac{\delta^2}{\delta \theta^2} \log(f_\theta(x))\right] $$ 
		\end{enumerate}
	\end{remark}

	\begin{exercise}
		Let $\underline{X} \sim Poi(\theta)$, Y=$\sum x_i$, and Y $\sim Poi(n\theta)$. What is $I(\theta)$?

		\begin{align*}
			I(\theta) &= E\left((\frac{\delta}{\delta \theta}\log f_\theta(\underline{X})^2) \right)\\
					&=E\left((\frac{\delta}{\delta \theta}\log \prod \frac{\theta^x e^{-\theta}}{x!})^2) \right)\\
					&= E\left((\frac{\delta}{\delta \theta}\log \frac{\theta^{\sum x_i} e^{-\theta n}}{\sum x_i!})^2 \right)\\
					&= E\left(\frac{\delta}{\delta \theta}(-n\theta + \sum x_i \log \theta - \sum \log x_i!)^2  \right)\\
					&= E\left((-n+\frac{\sum x_i}{\theta})^2 \right)\\
					&= E\left( n^2 - 2(\frac{\sum x_i}{\theta})(n)+(\frac{\sum x_i}{\theta})^2 \right)\\
					&= n^2 -\frac{2n * E(\sum x_i)}{\theta} + \frac{E((\sum x_i)^2)}{\theta^2}\\
					\\
					&=n^2 - \frac{2n*E(Y)}{\theta} + \frac{E(Y)^2}{\theta}\\
					&=\frac{n}{\theta}
		\end{align*}
	\end{exercise}

	Note: I($\theta$) is equal to information in the whole sample, but sometimes it's just one sample (based on context).

	If we assume (as in any exponential family):

	$$E_\theta \left(\frac{\delta^2}{\delta \theta^2}\log f_\theta(\underline{x})\right) = \frac{\delta^2}{\delta \theta^2}\int{\log f_{\theta}(\underline{x})}f_\theta(\underline{x})dx$$
	  
	then, the observed information is 
	$$I(\theta) = - E_\theta \left(\frac{\delta^2}{\delta \theta^2}\log f_\theta(\underline{x})\right) $$

	\begin{example}
		Let $X_i \dots X_n \stackrel{iid}{\sim} N(\mu,\sigma^2)$
	Find the information based on $\sigma^2$.

		Write 
	\begin{align*}
		\ell(\mu,\sigma^2) &=\log\left(\prod (2\pi\sigma^2)^{\frac{1}{2}} \exp\left\{\frac{-(x-\mu)^2}{2\sigma^2}\right\}\right)\\
		&= \frac{-n}{2} \log(2\pi\sigma^2) +\sum \left(\frac{-(x_i-\mu)^2}{2\sigma^2} \right)
		\end{align*}

	If we try to use the expected information: 

	$I(\sigma^2) = E\left(\frac{\delta}{\delta \theta}\left(\frac{-n}{2} \log(2\pi\sigma^2) +\sum \left(\frac{-(x_i-\mu)^2}{2\sigma^2}\right)\right)^2 \right)$
	which is a mess. 

	However, using the observed information:

	\begin{align*}
	I(\sigma^2) &= \frac{-n}{2\sigma^4} + \frac{1}{\sigma^6} E\left(\sum(X_i-\mu)^2\right)\\
		&= \frac{-n}{2\sigma^4}+\frac{n\sigma^2}{\sigma^6}\\
		&= \frac{n}{2\sigma^4}
	\end{align*}

	\end{example}

	\begin{example}
		Continued from previous example. 

		Define $S^2 = \frac{1}{n-1}\sum (X_i - \bar{X})^2$

		Note: $\frac{n-1}{\sigma^2}S^2 \sim \chi^2_{n-1}$

		Does $S^2$ acheive the C-R lower bound?

		\begin{align*}
		Var(S^2) &\ge \frac{\Psi^\prime(\theta)}{I(\theta)}\\
			\frac{2\sigma^4}{n-1} &\ge \frac{1}{I(\theta)}\\
			&\ge \frac{2\sigma^4}{n}
		\end{align*}
	\end{example}

	\begin{remark}
		Reread 7.3 and 6.2 in Casella and Berger
	\end{remark}

	What would give equality? When W($\underline{X}$) is a linear function of $\frac{\delta}{\delta \theta} \log f_\theta(\underline{X})$ which leads to... 

	\begin{corollary}[Attainment]
		Let $X_1,\dots,X_n$ be iid f(x|$\theta$), where f(x|$\theta$) satisfies the conditions of the Cramer-Rao Theorem. Let $L(\theta|x) = \prod f(x_i|\theta)$ denote the likelihood function. If $W(X) = W(X_1,\dots,X_n)$ is any unbiased estimator of $\tau(\theta)$, then W(X) attains the Cramer-Rao Lower Bound iff
		$$ a(\theta)\left(W(x) - \tau(\theta)\right) = \frac{\delta}{\delta \theta} \log L(\theta|x) $$
		for some function a($\theta$).
	\end{corollary}

	\begin{definition}[Efficient]
		Any estimator that achievies the Cramer-Rao bound is called an \textbf{efficient estimator}.
	\end{definition}

	\begin{example}
		$X_1, \dots, X_n \stackrel{iid}{\sim}	N(\mu, \sigma^s)$ where both are unknown. Let $\hat{\mu}_{ML} = \bar{X}$ and $\bar{\sigma^2}_{ML} =\frac{1}{n} \sum(x_i - \bar{X})^2$. Are these estimators efficient?
		
			\begin{align*}
				Var(\mu_{ML}) &\ge \frac{1}{I(\mu)}\\
			 	&= \frac{\sigma^2}{n}\\
			 	I(\mu) &= nI_1(\mu)\\
			 		&= nE\left[\left(\frac{\delta}{\delta \mu}\log f_\theta(x_i)\right)^2\right]\\
			 		&= nE\left[\left(\frac{\delta}{\delta \mu} \frac{-1}{2}\log\left(2\pi\sigma^2\right) - \frac{(x_1 - \mu)^2}{2\sigma^2}\right)^2\right]\\
			 		&= n E\left[\left(\frac{(x_1 - \mu)}{\sigma^2}\right)^2\right]\\
			 		&= \frac{n}{\sigma^4} E(x_1^2 -2x_1\mu + \mu^2)\\
			 		&= \frac{n}{\sigma^4}(\sigma^2 + \mu^2 - 2\mu^2 + \mu^2)\\
			 		&= \frac{n}{\sigma^2}
				\end{align*}		
			In this case, $\text{Var}(\mu_{ML}) = \frac{1}{I(\mu)}$, thus the sample mean is an efficient estimator.\\

			For $\sigma^2$, 

			\begin{align*}
			\text{Var}(\hat{\sigma^2}_{ML}) &= \text{Var}(\frac{1}{n} \sum(x_i - \bar{X})^2)\\
				&= \frac{2(n-1)\sigma^2}{n^2}\\
				I_1(\theta = \sigma^2) &= E \left[ \left( \frac{\delta}{\delta \theta} \frac{-1}{2}\log \left(2\pi\theta\right) - \frac{(x_1 - \mu)^2}{2\theta}\right)^2\right]\\
				&= E\left[\left(\frac{-1}{2\theta} + \frac{(x_1 - \mu)^2}{2\theta^2}\right)^2\right]\\
				&= \text{Var}\left[\frac{-1}{2\theta} + \frac{(x_1 - \mu)^2}{2\theta^2}\right]\\
				&= \text{Var}\left[\frac{(x_1 - \mu)^2}{2\theta^2}\right]\\
				&= \frac{1}{4\theta^4} \text{Var}\left[(x_1 - \mu)^2\right]\\
				&= \frac{1}{4\theta^4} \left[E(x_1 - \mu)^4 - \theta^2\right]\\ 
				&= \frac{3\theta^2 - \theta^2}{4 \theta^4}\\
				&= \frac{1}{2\sigma^4}\\
				\left[\Psi^\prime(\theta)\right]^2 &= \left[E \left( \frac{1}{n}\sum(x_i - \bar{x})  \right) \right]^2\\
					&= \frac{(n-1)\theta}{n}	
			\end{align*}
			So this estimator does not achieve the lower bound and thus is not efficient.

	\end{example}

	\textbf{Issues with CR Bound:} \begin{enumerate}[label = (\roman*)]
		\item The C-R bound need not be attained by our favorite estimators.
		\item Our assumptions do not always hold, and hense can not be used to apply the theorem. Especially when the range depends on a parameter. 
	\end{enumerate}

	For an alternative approach to finding the UMVUE, see the Rao-Blackwell and Lehmann
	-Sheffe Theorems. 

\subsection{Loss Function Optimality}\index{Loss Function Optimality}
	
	\begin{definition}[Loss] L($\theta$, W($\underline{x}$)) assigns a nonnegative real value called the \textbf{loss} to our decision to estimate $\theta$ by W($\underline{X}$). General context: Decision Theory. 
	\end{definition}

	Typically L($\theta$,$\theta$) = 0 because nothing is lost if your decision is exactly correct. 

	\begin{example}
		\begin{eqnarray*}
			L(\theta, W(\underline{X})) &=& (\theta-W(X)^2) \text{\indent square error loss}\\
	\\
			&=& | \theta -W(X)| \text{   \indent absolute error loss}\\
	\\
			&=& \frac{W(X)}{\theta}-1-\log \frac{W(X)}{\theta} \text{  \indent Stein's loss}
		\end{eqnarray*}
	\end{example}

	\begin{definition}[Risk] Average loss for a given estimator. \textbf{Risk} of estimating $\theta$ by W($\underline{X}$) is
		\[
		R(\theta,W)=E\left(\theta, W(\underline{X} \right) )
		\]
		
	\end{definition}

	\textbf{Goal:} We want to construct $\delta$ such that risk is minimized for every $\theta \in \Theta$.

	Two Ways:
	\begin{enumerate}
		\item Restrict the space of estimators.
		\item Allow all possible estimators but use a weaker notation of optimality
			\begin{enumerate}
			 	\item Bayesian approach
			 	\item MiniMax approach 
			 \end{enumerate} 
	\end{enumerate}

	\begin{definition}[Bayes Risk]
		Recall we previously define risk as a function of a parameter and estimator - E(Loss). Now that $\theta$ can be random, we may define \textbf{Bayes risk} as the average risk.
		$$r_\pi(\delta) = \int R(\theta,\delta)\pi(\theta) d\theta = E(R(\theta,\delta)) = E(L(\theta,\delta(\underline{X}))) $$
	\end{definition}

	\begin{definition}[Bayes Estimator]
		For a give loss function, we define the \textbf{Bayes estimator} as the minimizer of the Rayes Risk. 
	\end{definition}
	Note: if $\delta^*$ minimizes $E(L(\theta,\delta(\underline{X}))|\underline{X}=x)$ for (almost) all x then $\delta^*$ is the Bayes Estimator. 

	We can conclude that under square error loss, the posterior mean is the Bayes Estimator.\\
	Wrap up: How to choose prior? Even choosing family of distribution is hard. Considering restraints and conjugacy can help.
	\begin{itemize}
		\item Use data to estimate prio parameters
		\item Assign priors on our priors "hyper priors". Priorception.
		\item Use a Jeffery's prior: reparameterize
	\end{itemize}

	\begin{exercise}
	If $X_1, \dots, X_n$ are iid with mean $\mu$ and varience $\sigma^2$ what is
	$$
	E \left(\sum_{i=1}^n (X_i-\bar{X})^2 \right)?
	$$

	Note:
	\begin{enumerate}
			\item $\sum{i=i}^n \left (X_i - \bar{X})^2 \right) = \sum (X_i^2) - n\bar{X}$
			\item Var(X) = $E(X^2)-E(X)^2$
			\item Var$(\frac{X_i}{n})=\frac{\sigma^2}{n^2}$
			\item Var$(\bar{X})=\frac{\sigma^2}{n}$
		\end{enumerate}	

	Thus,
	$$E (\sum X_i^2) = (\sigma^2 + \mu^2)n$$
	$$-n E(\bar{X}^2)= -n(\mu^2+\frac{\sigma^2}{n}) $$
	We may conclude, $\frac{1}{n-1}\sum(X_i-\bar{X})^2$ is an unbiased estimator of $\sigma^2$ called $S^2$. 
	\end{exercise}

\begin{definition}[Minimax Risk]
	The minimum risk of an estimater $\delta$:
			$$\inf_{all \delta} \sup_{\theta \in \Theta} R(\delta, \theta) $$
	The Minimax risk is the minimum of the max risks. We say $\delta^*$ is minimax if
	$$\sup_{\theta \in \Theta} R(\delta^*, \theta) = \inf_{all \delta} \sup_{\theta \in \Theta} R(\delta, \theta)$$
	\end{definition}	

\begin{theorem}
	The Bayes risk, $r_\pi(\delta)$, satistifes 
	$$r_\pi(\delta) = \int R_\pi(\delta|\underline{x})m(\underline{x})d\underline{x}$$
	where $m(\underline{x}) = \int f_\theta(\underline{x}) \pi(\theta)d\theta$ and $R_\pi(\delta|\underline{x}) = \int l(\delta, \theta) \pi(\theta|\underline{x})d\theta$ is called the posterior risk. If $\delta^*(\underline{x})$ minimizes the risk, then it is a Bayes estimator 
\end{theorem}

\begin{proof}
\begin{align*}
	r_\pi(\delta) &= \int_\Theta R(\theta,\delta) \pi(\theta)d\theta\\
	&= \int_\Theta \int_{x^n} l(\theta,\delta) f_\theta(\underline{x})d\underline{x } \pi(\theta)d\theta\\
	\pi(\theta|\underline{x}) &= \frac{f_\theta (\underline{x})\pi(\theta)}{\int f_\theta (\underline{x}) \pi(\theta)d\theta}\\
	&= \frac{f_\theta(\underline{x})\pi(\theta)}{m(x)}
\end{align*}
	
\end{proof}

\begin{theorem}
	Suppose:
	\begin{enumerate}[label = \alph*)]
		\item $L(\theta, \delta) = (\theta - \delta)^2$. Then the Bayes estimator is the posterior mean.
		\item $L(\theta, \delta) = |\theta - \delta|$. Then the Bayes estimator is the posterior median.
		\item $L(\theta, \delta) = \begin{cases}
			0, & |\theta-\delta| \le \epsilon\\
			1, & |\theta-\delta| \> \epsilon\\
		\end{cases}$, then the Bayes estimator is the midpoint of an interval of length 2$\epsilon$ where $T(\theta|\underline{x})$ is maximized. 
		\item $L(\theta, \delta) = \begin{cases}
			0, & \theta=\delta|\\
			1, & \theta \neq \delta\\
		\end{cases}$, then the Bayes estimator is the maximum aposterior.
	\end{enumerate}
\end{theorem}

\subsection{Methods of Generating Reasonable Estimators}\index{Methods of Generating Reasonable Estimators}

Epmirical Bayes - parameters of the prior are estimated from data.

Hierarchichal Bayes - we impose a hierarch of priors on the parameters of $\theta$.

\textbf{Choice of Prior:} We want the prior impose \textbf{non-informative prior} where each value of the parameter is equiprobably. However, these priors could be \textbf{improper}, i.e. they don't integrate to 1 (in fact, not integrable at all). 

\begin{definition}[Jeffreys Prior]
	$\pi(\theta) \propto \sqrt{I(\theta)}$ If $\int \sqrt{I(\theta)} d\theta < \infty$ then $\pi$ is a proper prior on $\theta$. Else, we treat it as an improper prior, which is often the case for Jeffreys priors.
	\begin{remark}
		\begin{enumerate}
			\item Jeffreys priors works well for single parameter models ,but has some issues while dealing with multiple parameters.
			\item In general, Jeffreys prior is not a conjugate prior, but can be obtained as a limit of a conjugate prior family.
		\end{enumerate}
	\end{remark}
\end{definition}

\begin{definition}[Least Favorable Prior]
	A prior, $\pi$, is \textbf{least favorable} if
	$$r_\pi(\delta^*) \ge r_\pi(\delta^*)$$ for all priors, $\pi^\prime$
\end{definition}

\begin{theorem}
	Suppose $\pi$ is a prior on $\Theta$ such that
	\begin{align*}
		r_\pi(\delta^*) &= \int R(\theta, \delta^*) \pi(\theta) d\theta\\
		&= \sup_\theta R(\theta, \delta^*) (*)
	\end{align*}
	Then, 
	\begin{enumerate}[label = \alph*)]
		\item $\delta^*$ is minimax
		\item $\pi$ is the lease favorable prior.
		\item If $\delta^*$ is unique Bayes Estimator, then it is also unique minimax estimator.	
	\end{enumerate}
\end{theorem}

\textbf{Question:} When is (*) satisfied?
\begin{enumerate}
	\item $R(\theta, \delta^*)$ is constant with respect to $\theta$ (i.e. does not depend on $\theta$).
	\item $\Omega = {R(\theta, \delta^*) = \sup_{\theta^prime} R(\theta^\prime, \delta^*)}$. If $\pi(\Omega) = 1$ then (*) is satisfied.
\end{enumerate}

\begin{example}
	$X_1, \dots, X_n \stackrel{iid}{\sim} N(\mu, \sigma^2)$ where $\sigma^2$ is known. Is $\bar{X}$ minimax under squared error loss?

	 Let $\pi (\mu) \sim N(\tau, b^2)$

	$\delta^*$ is the Bayes estimator given by the posterior mean. 

	$\delta^*(\underline{x}) = \frac{\frac{n \bar{x}}{\sigma^2} + \frac{\tau}{b^2}}{\frac{n}{\sigma^2} + \frac{1}{b^2}}$

	As $b \rightarrow \infty$, $\delta^*(\underline{x} \rightarrow \bar{X}$. ($\pi$ is an improper prior for $\mu$ implies it's the Jeffreys prior for $\mu$).

	\begin{align*}
		R(\mu, \delta^*) &= E[L(\mu, \delta^*)]\\
		&= E[(\mu - \delta^*)^2]\\
		&= E[(\alpha \bar{x} + (1-\alpha)\tau - \mu)^2]\\
		&=E[\alpha^2(\bar{x}-\mu)^2 + (1-\alpha)^2(\tau - \mu)^2 + 2\alpha(\bar{x}-\mu)(1-\alpha)(\tau - \mu)]\\
		&= \alpha^2*E(\bar{x} - \mu)^2 + (1-\alpha)^2(\tau - \mu)^2\\
		&= \alpha^2 \text{Var}(\bar{x}) + (1-\alpha)^2(\tau - \mu)^2\\
		&= \frac{\alpha^2 \sigma^2}{n} + (1-\alpha)^2(\tau-\mu)^2\\
		r_\pi(\delta^*) &= \int R(\mu, \delta^*) \pi(\mu)d\mu (\text{Note: Risk from line above.})\\
		&= \int [\frac{\alpha^2\sigma^2}{n} + (1-\alpha)^2(\tau - \mu)^2]\pi(\mu)d\mu\\
		&= \frac{\alpha^2 \sigma^2}{n} + (1 - \alpha)^2 b^2\\
		\text{Substituting for $\alpha$, we get	}\\
		r_\pi(\delta^*) = \frac{1}{\frac{n}{\sigma^2}+\frac{1}{b^2}}\\
		&= \int R(\mu, \delta) \pi(\mu)d\mu\\
		&\le \sup_\mu R(\mu, \delta)\\
		\text{In other words},\\
		\sup_\mu R(\mu, \delta) \ge \frac{1}{\frac{n}{\sigma^2}+\frac{1}{b^2}}\\
		&\ge \sup_b \left(\frac{1}{\frac{n}{\sigma^2}+\frac{1}{b^2}}\right)\\
		&= \frac{\sigma^2}{n}, \forall \delta\\
		&\implies \inf_\delta \sup\mu R(\mu, \delta) = \frac{\sigma^2}{n} = R(\mu, \bar{x}) = \sup_\mu R(\mu, \bar{x})
	\end{align*}
	Therefore, $\bar{x}$ is minimax.
\end{example}

\begin{definition}[Admissibility]
An estimator $\delta$ is said to be inadmissable if there exists $\delta^\prime$ that dominates it, i.e. R($\theta, \delta) \ge R(\theta, \delta^\prime)$ for all $\theta$ and $R(\theta, \delta) > R(\theta, \delta^\prime)$ for some $\theta$. If no such $\delta^\prime$ exists, then $\delta$ is \textbf{admissible}. 
\begin{enumerate}[label = \alph*)]	
 		\item UMVUE need not be admissible as biased estimators can yield smaller MSE.
 		\item But the situation is different when dealing with minimax estimators.
 		\item If $\delta$ is a minimax estimator and it is dominated by $\delta^\prime$, then $\delta^\prime$ is also minimax and therefore we prefer $\delta^\prime$.
 		\item So it is important to ascertain whether a given minimax estimator is admissable or not.
 	\end{enumerate}	 	
\end{definition}

\begin{theorem}
	Any unique Bayes estimator is \textbf{admissible}. 
	$$\inf_\delta r_\pi(\delta)$$
\end{theorem}

\textbf{Proposition:} If an estimator is unique minimax, then it is admissible. 

\textbf{Proposition:} If $\delta$ has a constant risk and is admissible then it is minimax.

\section{Practice Problems}

Chapter 7 Exercises (pp. 355 to 367): 7.2, 7.6, 7.7, 7.10, 7.12, 7.13, 7.14, 7.19, 7.20, 7.21, 7.22, 7.23, 7.25, 7.30, 7.37, 7.38, 7.44, 7.47, 7.48, 7.49, 7.52\\
\\

	\begin{problem} (C\&B 7.2) Let $X_1, \dots, X_n$ be a random sample from a Gamma($\alpha, \beta$) population.

		\begin{enumerate}[label=(\alph*)] 
			\item Find the MLE of $\beta$, assuming $\alpha$ is known.
			\item If $\alpha$ and $\beta$ are both unknown, there is no explicit formula for the MLEs of $\alpha$ and $\beta$, but the maximum can be found numberically. The result in part (a) can be used to reduce the problem to the maximization of a univariate function. Find the MLEs for $\alpha$and $\beta$ for the data in Exercise 7.10(c).
		\end{enumerate}

			Solution.

			(a) Gamma distribution pdf: $\frac{\beta^\alpha}{\Gamma(\alpha)}x^{\alpha-1}e^{-\beta x}$\\
			\\
			$L(\beta) = \frac{\beta^{n\alpha}}{\Gamma(\alpha)^n}[\prod x]^{\alpha-1}e^{-\beta \sum x}$\\

			$\ell(\beta) = n\alpha\log\beta - n\log\Gamma(\alpha) +(\alpha-1)\log(\prod x_i) - \beta \sum x_i $\\

			$\frac{\delta \ell}{\delta \beta} = \frac{n\alpha}{\beta} - \sum x_i$\\
			\\
			$\hat{\beta} = \frac{n\alpha}{\sum x_i}$

			\begin{remark}
				The density of Gamma can also be in the form with $\frac{1}{\beta}$ which would mean $\hat{\beta}$ would be $\frac{\sum x_i}{n\alpha} $.
			\end{remark} 

			(b) The new likelihood with $\hat{\beta}$ plugged in is:
			$$ L(\beta) = \frac{(\frac{n\alpha}{\sum x_i})^{n\alpha}}{\Gamma(\alpha)^n}[\prod x]^{\alpha-1}e^{-(\frac{n\alpha}{\sum x_i}) \sum x}$$ 

			To solve...use a computer.\\
			\\

	\end{problem}
	
	\begin{problem} (C\&B 7.6) Let $X_1, \dots, X_n$ be a random sample from the pdf
	$$f(x|\theta) = \theta x^{-2}, 0<\theta\le x < \infty $$

	\begin{enumerate}[label=(\alph*)]
		\item What is a sufficient statistic for $\theta$?
		\item Find the MLE of $\theta$.
		\item Find the method of moments estimator of $\theta$?\\
		\\
	\end{enumerate}
		

		Solution. 

		(a) Use Factorization Theorem.\\ 
		$L(\theta) = \theta^n \prod(x_i^2) \prod I\{\theta \le x < \infty\}$ where $\prod I\{\theta \le x < \infty = x_{(1)}$\\
		Thus the sufficient statistic is the minimum x. \\
		\\
		(b) L($\theta$|x) = $\theta^n \prod(x_i^2) \prod I\{\theta \le x < \infty \}$. $\theta^n$ is increasing in $\theta$. The second term does not involve $\theta$. So to maximize $L(\theta|x)$, we want to make $\theta$ as large as possible. But because of the indicator function, $L(\theta|x) = 0 if \theta > x(1)$. Thus, $\hat{\theta} = x_{(1)}$.\\
		\\
		(c) $E(X) = \int^\infty_\theta \theta x^{-1} dx = \theta \log x|^\infty_\theta = \infty$. Thus the method of moments estimator of $\theta$ does not exist.\\
		\\

	\end{problem}

	\begin{problem} Let $X_1, \dots, X_n$ be iid with one of two pdfs. If $\theta$ = 0, then
	$$ f(x|\theta) = \begin{cases} 1,\text{ if } 0<x<1 \\ 0, \text{ otherwise} \end{cases}$$
	$$ f(x|\theta) = \begin{cases} \frac{1}{2\sqrt{x}},\text{ if } 0<x<1 \\ 0, \text{ otherwise} \end{cases}$$
	Find the MLE of $\theta$.\\
	\\

	Solution.\\
	\\
	$$ L(0|x) = \begin{cases} 1,\text{ if } 0<x_i<1 \\ 0, \text{ otherwise} \end{cases}$$
	while if $\theta=1$, then
	$$ L(1|x) = \begin{cases} \prod \frac{1}{2\sqrt{x}},\text{ if } 0<x_i<1 \\ 0, \text{ otherwise} \end{cases}$$

	Thus, 

	$$\hat{\theta}_{MLE} = \begin{cases} 1,\text{ if } \prod \frac{1}{2\sqrt{x}}\ \\ 0, \text{ otherwise} \end{cases} $$
		
	\end{problem}


%------------------------------------------------

% \section{Citation}\index{Citation}

% This statement requires citation \cite{book_key}; this one is more specific \cite[122]{article_key}.

% %------------------------------------------------

% \section{Lists}\index{Lists}

% Lists are useful to present information in a concise and/or ordered way\footnote{Footnote example...}.

% \subsection{Numbered List}\index{Lists!Numbered List}

% \begin{enumerate}
% \item The first item
% \item The second item
% \item The third item
% \end{enumerate}

% \subsection{Bullet Points}\index{Lists!Bullet Points}

% \begin{itemize}
% \item The first item
% \item The second item
% \item The third item
% \end{itemize}

% \subsection{Descriptions and Definitions}\index{Lists!Descriptions and Definitions}

% \begin{description}
% \item[Name] Description
% \item[Word] Definition
% \item[Comment] Elaboration
% \end{description}

%----------------------------------------------------------------------------------------
%	CHAPTER 2
%----------------------------------------------------------------------------------------

\chapter{Principles of Data Reduction}

\section{Sufficiency Principle}

\begin{definition}[Sufficient]
	T(X) is \textbf{sufficient} for $\theta$ if the distribution of X|T($\underline{X}$) does not depend on $\theta$. Note that sufficient statistics are not unique.
\end{definition}

\begin{theorem}[Factorization Theorem]
	If we have $\underline{X} \sim f_\theta(\underline{x})$ then T is sufficient if and only iff we can write $f_\theta$ as
	$$g(T(x),\theta)h(x)$$
	for some g and h.
\end{theorem}

\begin{example}
	$X_1, \dots, X_n \stackrel{iid}{\sim}$ Bern($\theta$)\\
	pdf: $\theta^x(1-\theta)^{1-x}$

	Joint pdf = $f_\theta(x) = \theta^{\sum x_i}(1-\theta)^{n-\sum x_i}$

	So, T($\underline{x}$) = $\sum x_i$
\end{example}


Q: What is the connection to Section 7.3?\\
A: \begin{align*}
Var_\theta(W(\underline{X})) &= Var\left[E(W(X)|T(X))\right] + E\left[Var(W(X)|T(X)) \right]\\
&\ge Var\left[E(W(X)|T(X))\right]	
\end{align*} 
Note: $E_\theta \left\{E_\theta[W(X)|T(X)] \right\} = E_\theta(W(X))$

$E_\theta[W(X)|T(X)]$ does not depend on $\theta$. So this is a legitimate estimator since T(X) is sufficient and its varience is smaller than any estimator with same mean. \\
\\
General Idea of Sufficiency: If T(X) is sufficient for $\theta$, then all information in X about $\theta$ is captured in T(X). \\
\\
More technically, conditioning T(X), the remaining randomness does not depend on $\theta$.

\begin{remark}
	Know Thm 6.2.6 - Factorization Theorem
\end{remark}

\begin{definition}[Minimal Sufficient]
	 T(X) is \textbf{minimal sufficient} if for any other sufficient statistics, $T^\prime(X)$, $T(X)$ is a function of $T^\prime(X)$.
	 Note: "T(X) is a function S(X)" means that S(x) = S(y) implies T(x) = T(y).
\end{definition}

\begin{example}
	 $X_1, \dots, X_n \stackrel{iid}{\sim} N(\mu,\sigma^2)$. Find possible T($\mu,\sigma^2$).\\
	 \begin{align*}
	 f(\underline{x}) &= \prod (2\pi\sigma^2)^{\frac{1}{2}} e^{\frac{-(x-\mu)^2}{2\sigma^2}}\\
	 &= (2\pi\sigma^2)^{\frac{-n}{2}}\exp\left\{\frac{-\sum x_i^2 + 2\mu \sum x_i - n\mu^2}{2\sigma^2}\right\}
	 \end{align*}
Therefore our pair of sufficient statistics are $\sum x_i^2$ and $\sum x_i$. But are the minimal?
\end{example}

\begin{theorem} $T(\underline{x}) = T(\underline{y})$ if and only if $\frac{f_\theta (\underline{x})}{f_\theta(\underline{y})}$ does not depend on $\theta$, implying T($\underline{x}$) is minimal sufficient statistic. 
\end{theorem}

\begin{remark}
	Minimal sufficient statistics are not unique (any 1-1 transformation of a minimal sufficient statistics is also minimal sufficient). A sufficient statistic may have components that do not contain any information about $\theta$. Such components are called \textbf{ancillary statistics}.
\end{remark}

\section{Exponential Families}

	\begin{definition}{Exponential Family} An \textbf{exponential family} of densities is the set of densities given by 
	$$f_{\underline{\theta}}(\underline{x})= h(\underline{x})c(\underline{\theta})\exp\left\{\sum w_i(\underline{\theta})T_i(\underline{x})\right\} $$
	As $\theta$ takes values in some set $\Theta$, this equation may be rewritten as
	$$f_{\underline{\eta}}(\underline{x})= \exp\left\{\textcolor{red}{\underline{\eta}}^T\textcolor{green}{\underline{T}(\underline{x})}-\textcolor{blue}{A(\underline{\eta})}\right\}\textcolor{magenta}{h(\underline{x})} I_B(\underline{x}) $$
	The latter equation is called the \textbf{canonical/natural parameterization}.
	\end{definition}

	\begin{remark}
		The $T_i$/$T(x)$ terms are sufficient statistics.
	\end{remark}

	\begin{example}
		Deriving the canonical parameterization for a Normal distribution.
		\begin{align*}
			f_{(\mu,\sigma^2)} &= k \frac{1}{\sqrt{2\sigma^2}}\exp\left\{\frac{-1}{2\sigma^2}(x-\mu)^2 \right\}\\
			&= k \exp\left\{\textcolor{red}{\frac{1}{2\sigma^2}} \textcolor{green}{(-x^2)} + \textcolor{red}{\frac{\mu}{\sigma^2}} \textcolor{green}{(x)} - \textcolor{blue}{\frac{\mu^2}{2\sigma^2}+\frac{1}{2}\log((2\sigma^2))} \right\}\\
			&= k \exp\left\{\textcolor{red}{\eta_1} \textcolor{green}{(-x^2)} + \textcolor{red}{\eta_2} \textcolor{green}{(x)} - \textcolor{blue}{\frac{\eta_2^2}{4\eta_1} - \frac{1}{2}\log\eta_1 } \right\}
		\end{align*}

		$\eta_1 = \frac{1}{2\sigma^2} \Leftrightarrow \sigma^2 = \frac{1}{2\eta_1}$\\
		$\eta_2 = \frac{\mu}{2\sigma^2} \Leftrightarrow \mu = \frac{\eta_2}{2\eta_1}$\\
		$A(\underline{\eta})= \frac{\eta_2^2}{4\eta_1} - \frac{1}{2}\log\eta_1 $

	\end{example}

	Cool fact: The partial derivatives of A($\eta$) gives the expectations of T. The second partials give the covarience of T.

	Observe $T(x)=(-x^2,x)$

		\begin{align*}
			E(-x^2) &= \frac{\delta}{\delta\eta_1}A(\underline\eta)\\
				&= \frac{-\eta_2^2}{4\eta_1^2}-\frac{1}{2\eta_1}\\
				&= \frac{-\mu^2 * 4\sigma^4}{4 \sigma^4} - \frac{2\sigma^2}{2} \\
				&=-\mu^2-\sigma^2
		\end{align*}

		\begin{align*}
			E(x) &= \frac{\delta}{\delta\eta_2}A(\underline{\eta})\\
				&= \frac{\eta_2}{2\eta_1}\\
				&= \mu
		\end{align*}

		\begin{align*}
			Var(x^2) &= \frac{\delta^2}{\delta^2 \eta_1}\\
				&= \frac{\eta_2^2}{2\eta^3_1}+\frac{1}{2\eta_2}\\
				&= 4\mu^2\sigma^2+2\sigma^4
		\end{align*}

		\subsection{Distribution Families}
\begin{description}
	
		\item[Poisson] 
		
		\begin{align*}
		f_\lambda &= e^{-\lambda}\frac{\lambda^x }{x!} I\{x \in \mathbb{N} \} \\
			&= \exp \left\{\textcolor{red}{(\log\lambda)} \textcolor{green}{x}-\textcolor{blue}{\lambda} \right\}\\
		\eta &= \log\lambda\\
		T(\underline{x}) &= x\\
		A(\eta) &= \lambda = e^{\log\lambda}=e^\eta\\	
		\end{align*}


		\item[Beta] 
		\begin{align*}
		f_{\alpha,\beta}(x) &= \frac{\Gamma(\alpha+\beta)}{\Gamma(\alpha)+\Gamma(\beta)} x^{\alpha-1} (1-x)^{\beta-1} I\left\{0<1x<1 \right\}\\
			&=  \exp \left\{\log \Gamma(\alpha+\beta)-\log\Gamma(\alpha)-\log\Gamma(\beta)+(\alpha-1)\log x+(\beta-1)\log(1-x) I\{0<x<1\} \right\}\\
			\eta_1 &= \alpha-1\\
			\eta_2 &= \beta-1\\
			T_1(x) &= \log(x)\\
			T_2(x) &= \log(1-x)-x\beta+\alpha\log(\beta)-\log(\Gamma(\alpha))
		\end{align*}

		\item[Gamma] 
		\begin{align*}
			f_{\alpha,\beta}(x) &= \frac{\beta^\alpha x^{\alpha-1}}{\Gamma(\alpha)}\exp\left\{-x\beta \right\} I\{x>0\}\\
			&= \exp \left\{(\alpha-1)\log(x)-\beta(x) + (\eta_1+1)\log(-\eta_2)+\log(\Gamma(\eta_1+1))\right\}\\
			\eta_1 &= \alpha-1\\
			\eta_2 &= \beta\\
			T_1(x) &= \log(x)\\
			T_2(x) &= x\\
			A(\eta) &= (\eta_1+1)\log(-\eta_2)+\log(\Gamma(\eta_1+1))\\
		\end{align*} 
		
		\item[Binomial]
		\begin{align*}
			f_p(x) &= \frac{n!}{x!(n-x)!}p^x(1-p)^{n-x}I\{0\dots,n\} \\
				&= \begin{pmatrix}  n \\  x  \end{pmatrix} \exp\left\{x(\log\frac{p}{1-p})-(-n\log(1-p)) \right\}\\
			\eta&=\log\frac{p}{1-p}\\
			T(x)&= x\\
			A(\eta)&= -n\log(1-p)
		\end{align*}

\end{description}

Why is E(T(x)) = $\frac{\delta}{\delta\eta}A(\eta)$?\\
Notice: $e^{A(\eta)}$ must be the integral of $\exp\left\{\textcolor{red}{\underline{\eta}}^T\textcolor{green}{\underline{T}(\underline{x})}-\textcolor{blue}{A(\underline{\eta})}\right\}\textcolor{magenta}{h(\underline{x})}$.\\
Differentiating on both sides and using the fact that for an expoonential family we may switch integral and derivative, we obtain:
\begin{align}
	\frac{\delta}{\delta\eta}A(\eta)e^{A(\eta)}&= e^{A(\eta)}\int\exp\left\{\eta^TT(x) \right\}h(x)[T(x)]e^{-A(\eta)}\\
	&= E(T(x))\\
\end{align}

Recall, $f_\eta (x)$ must always integrate to 1!

\begin{remark}
	Read Casella \& Berger Chapter 7.2
\end{remark}

\section{Complete Statistic}
Recall, minimal statistics contain no redundant or unnessicary information. What might redundant information look like? From T(x) we may construct $g_1$ and $g_2$ such that $E_\theta(g_1(T))=E_\theta(g_2(T))$ for all $\theta$ which we want to avoid.

\begin{definition}[Complete Statistic] (T) is obtained only if $g=0$ gives $E_\theta(g(T))=0$ for all $\theta$. 
	
\end{definition}

\begin{theorem}Complete and sufficient inplies minimal statistics.
	
\end{theorem}

\begin{theorem} If $f_\theta(\underline{x}) \exp\left\{\eta^T T(x)=A(\eta)\right\} h(x)$ is a cononical exponential family of full rank then T is a complete and sufficient statistic.
	
\end{theorem}

\begin{definition}{Full Rank} \textbf{Full rank} means
	\begin{itemize}
		\item Parameter space H contains an open set
		\item T does not satisfy any linear constraint (i.e. the T are linearly independent).
	\end{itemize}
\end{definition}

\begin{definition}[Open Set] Let k be dimention of $\eta$. To say that H contains an \textbf{open set} means there exists $\begin{cases} \epsilon >0 \\ \eta_0 \in H \end{cases}$ such that $B(\eta_0;\epsilon)\subset H$
	
\end{definition}

\begin{theorem}[Basu's Theorem] Any complete and sufficient statistic (minimum sufficient) is independent of any ancillary statistic (that is, any statistic whose distribution does not depend on $\theta$). 
	
\end{theorem}

\subsection{Sufficiency and Unbiasness}

Suppose $E_\theta(W(X)) = \theta$ for all $\theta$, thus W(X) is unbiased. Then Var[$E(W(X)|T)]\le Var(W(X))$. But what if T is not sufficient? Wel...then you won't get an estimator.

\begin{example} Suppose that $X \sim Unif(0,\theta)$. Find the UMVUE of $\sin\theta$.\\
\\
Solution. \\
pdf: $f(x)=\frac{1}{b-a} = \frac{1}{\theta}$\\
E(X)= $\frac{\theta}{2}$\\
Thus 2X is unbiased for $\theta$. But is $\sin(2X)$  unbiased for $\sin(\theta)$?\\
NO!\\

What about cos(2X)?
\begin{align*}
E(\cos(2X))&= \int^\theta_1 \cos(2X)\frac{1}{\theta}dx\\
	&= \frac{1}{\theta}\left(\frac{\sin(2X)}{2} \right)|^\theta_0\\
	&=\frac{\sin(2\theta)}{2\theta}
\end{align*}
No! Blarg. Let's bust out some calculus.

\begin{align*}
	\int^\theta_0 \textbf{estimator(x)} dx &= \theta \sin(\theta)\\
		\textbf{estimator} &= \sin\theta+\theta\cos(\theta)\\
			&= x\sin(x)+x\cos(x)
\end{align*}
 	
 \end{example} 

 \begin{theorem}[Rao-Blackwell] Let W(X) be an extimator with $E_\theta(W(x))=\tau(\theta)$. If T(X) is a sufficient statistic, then an alternative, unbiased estimator of $\tau(\theta)$ whose varience is uniformly not worse than that of W(X) is
 	$$\phi(X)=E_\theta(W(X)|T(X))$$.
 \end{theorem}

 \begin{proof}
  	Based on conditional arguments. 
 $$E(W)=E(E(W(X)|T(X)))$$ Var(W)=Var(Cond Exp)+ E(Cond Var)
  \end{proof} 

 \begin{theorem}
 	If W(X) is UMVUE for $\tau(\theta)$ then it is the unquie UMVUE, aka the best unbiased estimator for all $\theta$. 
 \end{theorem}

 \begin{theorem}{7.3.20}
 	If $E_\theta(W(X))=\tau(\theta)$, then W(X) is UMVUE for $\tau(\theta)$ iff W(X) is uncorrelated with any S(X) such that $E_\theta(S(X))=0$. 
 \end{theorem}

 \begin{theorem}[Lehmann-Scheffe Theorem]
 	If T(X) is a complete and sufficient statistic, then $\phi(T)$ is the unique UMVUE of its expectation (see Rao-Blackwell).
 \end{theorem}

\begin{definition}
	A funciton $\phi: \mathbb{R}^d \rightarrow \mathbb{R}$ is convex if domain ($\phi$) is a convex set and if for all $x, y \in dom(\phi)$ and $0 \le \alpha \le 1, \phi(dx + (1-\alpha)u) \le \alpha \phi(x)+(1-\alpha)\phi(y).$ If strict equality holds for all $x\neq y$, then $\phi$ is strictly convex.
\end{definition}

\textbf{Characterization of Convexity:}
\begin{enumerate}
	\item First-order conditions: suppose $\phi$ is differentiable (i.e. the gradient $\nabla\phi$ exists at each point in dom($\phi$) which is open). Then $\phi$ is convex iff dom($\phi$) is convex and $\phi(y)\ge \phi(x) + \nabla\phi^T(x)(y-x), \forall x, y.$
	\item Second-order conditions: suppose $f$ is twice differentiable. That is, the Hessian $\nabla^2\phi$ exists at each point in dom($\phi$), which is open. Then $\phi$ is convex iff dom($\phi$) is convex and $\nabla^2\phi$ is positive semidefinite.
\end{enumerate}

 \begin{definition}[Jenson's Inequality]
 If h(x) is a convex function then 
 $$h(E(x))\le E(h(x)) $$
for any random variable x.  	
 \end{definition}

\begin{theorem}
	If $\phi$ is convex, then $\phi(E(X)) \le E(\phi(X))$ and the strict inequality holds if $\phi$ is strictly convex.
\end{theorem}

\begin{corollary}
	Let X be a idscrete random variable with $P(X=x_i) = \alpha_i, i - 1,...l$. Apply the above theorem. 
\end{corollary}

\textbf{Uniqueness of $\delta^*$:} 
$$\delta^* \in \arg\inf_\delta r_\pi(\delta)$$
We are done if $r_\pi(\delta)$ is strictly convex for $0<\alpha<1$, consider 
\begin{align*}
	r_\pi(\alpha \delta + (1-\alpha)\delta_2) &= \int R(\theta, \alpha \delta_1 + (1-\alpha)\delta_2) \pi(\theta) d \theta\\
	&= \int E\left[L(\theta, \alpha \delta(1 + (1-\alpha)\delta_2)))\right]\pi(\theta) d \theta\\
	&< E\left[\alpha L(\theta, \delta) + (1-\alpha)L(\theta, \delta_2)\right]\pi(\theta) d \theta\\
	 &= \int(\alpha E(L(\theta, \delta_1)+(1-\alpha)E(L(\theta, \delta_2)))))\\
	 &= \alpha r_\pi(\delta_1) + (1-\alpha)r_pi(\delta_2)
\end{align*}	
Thus, $r_\pi$is strictly convex, which implies that $d^*$ is the unique minimizer. 

\textbf{Generalized Version of Rao-Blackwell Theorem}
Let w be an estimator of $g(\theta)$ and T be a sufficient statistic. Let $L(\theta,\delta)$ be a strictly convex function in its second argument.

Suppose $E(W)<\infty$ and $R(\theta, W) = E(L(\theta,W))<\infty$. Define $\eta(T) = E(W|T)$. Then $R(\theta, \eta(T)) < R(\theta, W)$ unless $\eta(T)=W$ with probability 1.

\begin{theorem}
	Suppose T is complete and sufficient. Then for every $g(\theta)$ that can be estimated by an unbiased estimator, there is one and only one unbiased estimator that is a function of T. In addition, this estimator uniformly minimizers the risk for any loss function $L(\theta, \delta)$ that is stictly convex in $\delta$. In particular, this estimator is UMVUE.
\end{theorem}

We say $\hat{\theta}$ is \textbf{efficient} if it achieves equality in the information bound. We say $\hat\theta_{ML}$ is \textbf{asymptotically efficient} if as $n\rightarrow \infty$, Var($\hat\theta_{ML}) \rightarrow $ the CR bound. 

\section{Major/Minorization Mini/Maximization}

Every EM Algorhythm is special case of a MiMax.

EM = Expectation Maximization

\begin{remark}
	Read 7.2.4 in Casella \& Berger
\end{remark}

\subsection{Bharath's Notes}

\textbf{Minoroization - Maximization}

\begin{enumerate}
	\item Constructing a minorization function. 
	\begin{align*}
		f(x) &\ge g(x,y), \forall x,y \in \Omega\\
		f(x) &= g(x,x), \forall x \in \Omega
	\end{align*}
	\item $X^{(t+1)} \in \text{argmax}_{x \in \Omega \subset \mathbb{R}^d } g(x,x^{(t)})$
	\item Return to step 1.
	\end{enumerate}

\textbf{Claim:} $f(x^{t+1}) \ge f(x^{(t)})$

\begin{proof}
	\begin{align*}
		f(x^{t+1}) &\ge g(x^{(t+1)}, x^{(t)})\\
		&\ge g(x^{(t)}, x^{(t)})\\
		= f(x^{(t)})
	\end{align*}
\end{proof}

\textbf{EM Algorithm:}
Goal: Maximize $L(\theta) = f_\theta(x_1, \dots, x_n)$ over $\theta \in \Theta \subset \mathbb{R}^d$.

\begin{definition}
	The Kullbach-Liebler divergence between two densities p and q is defined as 
	\begin{align*}
		KL(p||q) &= E_{x \sim q}(\log \frac{q(x)}{p(x)})\\
		&= \begin{cases}
			\int q(x)\log\frac{q(x)}{p(x)} dx, & if q<<p
		\end{cases}
	\end{align*}
	
\end{definition}

\subsection{Dave's Notes}

\begin{example}
	$X_1, \dots, X_n \stackrel{iid}{\sim} \text{Bern}(p), 0<p<1$
	\begin{align*}
		\hat{p}_{mm} &= \frac{1}{n} \sum \limits^n_{i=1} x_i\\
			&= p^{\sum x_i} (1-p)^{n - \sum x_i}
		L(p)\\  
		&= \prod\limits^n_{i=1} p^{x_i} (1-p)^{1-x_i}\\
		l(p) &= \sum x_i \log(p) + (n + \sum x_i) \log(1-p)\\
		\frac{d l(p)}{dp} &= \frac{\sum x_i}{p} + \frac{n - \sum x_i}{(1-p)}\\
			&\implies \hat{p}_{ML} = \frac{1}{n}\sum x_i\\
		\frac{d^2 l(p)}{dp^2} &= \frac{- \sum x_i}{p^2} - \frac{n - \sum x_i}{(1-p)^2} < 0
	\end{align*}
	and therefore $\hat{p}_{ML}$ maximizes L(p).
\end{example}

\subsection{How EM Algorhythms Work}
\begin{enumerate}[label = \alph*)]
	\item Choose a starting parameter, $\theta$. 
	\item [E-Step] Construct the (minimizing) function.
	\item Maximize this function of $\theta$. The maximizer will be next value of $\theta_0$. Call it $\theta_1$. 
	\item Return to (b) as long as we haven't converged.   
\end{enumerate}

Given $\mu_0$ we find after one EM iteration that
\begin{align*}
\mu_1 &=\frac{\sum Y_i}{n}+\mu_0(1-\frac{\mu}{n})\\
&= \mu_0(1-\frac{\mu}{n})+\frac{\sum Y_i}{\mu}(\frac{\mu}{n})\\
\end{align*}

Thus, $\mu_1$ always takes us $\frac{\mu}{n}$ of the way to the final answer at each iteration. This type of convergence is called \textbf{linear convergence}, which is a characteristic of EM algorhythems. The linear rate of convergence is governed by "amount of missingness" and is considered slow compared to other similar methods (e.g. Newton-Raphson is optimization method that enjoys quadradic convergence). However, EM algorhithem tend to trade more iterations for simpler iterations. 

Suppose we wish to maximize a fuction that is a product of sums or integrals. Taking the log gives a sum of logs of sums/integrals:
$$\sum^A\log\sum^B s_{ab}(\theta)$$ 
To simplify, ignore the sumation over A and take $f(\theta) = \log \sum s_{ab}(\theta)$. Fix a $\theta_0$. Define
$$W_{ob} \equiv \frac{S_b(\theta_0)}{\sum^c_{c_a} S(\theta_0)} $$ 

Goal: Maximize f($\theta)=\log \sum S_b(\theta)$ with fixed $\theta_0$.

Claim: Define $Q_0(\theta) \equiv \sum^B W_{ob} * \log S_b(\theta)$. \\
Then $f(\theta) -f(\theta_0) \ge Q_0(\theta) - Q_0(\theta_0)$. \\
\\
\\

To verify this claim, notice that Jenson's inequality says 
$$E(\log(\bullet)) \ge \log E(\bullet) $$

\begin{align*}
	Q_0(\theta)-Q_b(\theta_0)&= \sum^B W_{ob} \log(\frac{S_b(\theta)}{S_b(\theta_0)})\\
		&\le \log \sum^B W_{ob} \frac{S_b(\theta)}{S_b(\theta_0)}\\
		&\le \log \sum^B \frac{S_b(\theta_0)}{\sum S_c(\theta_0)} \frac{S_b(\theta)}{S_b(\theta_0)}\\
	f(\theta) - f(\theta_0) &= \log \frac{S_b(\theta)}{S_c(\theta_0)}	
\end{align*}


\begin{example}
	$X\sim f_\theta(x)=\lambda f_{\xi_1}(x)+(1-\lambda)f_{\xi_2}(x),  \theta=(\lambda, \xi) $
Intuition: To generate X acoording to mixture density flip coin with Heads probability of $\lambda$. If H, generate $X\sim f_{\xi_1}(x)$. If T, generage $X\sim f_{\xi_2} $. 

Write down the observed data likelihood, $\ell(\lambda, p_1, p_2)$. For a sample size, n, form a mixture of Binom(m, $p_1$) and Binom(m, $p_2$). 

$$\sum \log \left(\lambda \left(\begin{matrix} m \\ \xi \end{matrix}\right) p_1^{\xi_1}(1-p_1)^{m-\xi_1}+(1-\lambda)\left(\begin{matrix} m \\ \xi_2 \end{matrix}\right) p_2^{\xi_2}(1-p_2)^{m-\xi_2}\right) $$  
Now impliment EM algorithm.
\end{example}

Suppose that "complete data" consists of $x_{obs} \& x_{miss}$. 

Using previous notation, 
\begin{align*}
	S_b(\theta) &= P_\theta\left(X_{obs}=x_{obs}, X_{miss}=b \right)\\
	&\Rightarrow \sum^B S_b(\theta) = P_\theta\left(X_{obs}=x{obs} \right) \\
	\\
	W_{ob}&= \frac{P_{\theta_0}\left(X_{obs}=x_{obs}, X_{miss}=b \right)}{P_\theta\left(X_{obs}=x_{obs}\right)}\\ 
		&= P_{\theta_0}\left(X_{miss}=b|X_{obs}=x_{obs}\right)\\
		\\
	Q_0(\theta)&= E_{\theta_0}[\log P_\theta\left(X_{obs}=x_{obs}, X_{miss}=b\right)|X_{obs}=x_{obs}]	
\end{align*}


%----------------------------------------------------------------------------------------
%	CHAPTER 3
%----------------------------------------------------------------------------------------

\chapterimage{chapter_head_1.pdf} % Chapter heading image

\chapter{Hypothesis Testing}

\begin{definition}{hypothesis testing}
	We see that a hypothesis testing procedure is a rule that partitions the sampe space into we will accept (fail to reject) $H_0$ as true or not.
\end{definition}

Let's consider a simple null hypothesis, i.e. a hypothesis of the form:
$$H_0: \theta=\theta_0 $$
We will consider three different ways to determine whether $\theta_{true} = \theta_0$

\begin{enumerate}
	\item \textbf{Wald Test} $\hat{theta} - \theta_0$ - consider if this is large relative to the distribution. $\hat{\theta}$ should have if $H_0$ true. Recall, $\sqrt{n}(\hat{\theta}-\theta_0)\sim N(0,I^{-1}(\theta_0)$.
	\item \textbf{Likelihood Ratio Test} Consider $\ell(\hat{\theta}) - \ell(\theta_0)$ and determine whether this is large relative to $\chi^2$ distribution that it will approximately have under $H_0$. 
	\item \textbf{Rao Score Test} Consider $\nabla\ell(\theta_0)$ and determine if it is too far from 0 relative to it's true approximate distribution if $H_0$ true.  
\end{enumerate}

\begin{remark}
	Read 8.2.1 and 8.2.2
\end{remark}

\begin{theorem}
	If T(X) is sufficient for the family of distributions from which X is drawn and $\delta(X)$ is an unbiased estimater 
\end{theorem}
% \section{Table}\index{Table}

% \begin{table}[h]
% \centering
% \begin{tabular}{l l l}
% \toprule
% \textbf{Treatments} & \textbf{Response 1} & \textbf{Response 2}\\
% \midrule
% Treatment 1 & 0.0003262 & 0.562 \\
% Treatment 2 & 0.0015681 & 0.910 \\
% Treatment 3 & 0.0009271 & 0.296 \\
% \bottomrule
% \end{tabular}
% \caption{Table caption}
% \end{table}

% %------------------------------------------------

% \section{Figure}\index{Figure}

% \begin{figure}[h]
% \centering\includegraphics[scale=0.5]{placeholder}
% \caption{Figure caption}
% \end{figure}

%----------------------------------------------------------------------------------------
%	BIBLIOGRAPHY
%----------------------------------------------------------------------------------------

\chapter*{Bibliography}
\addcontentsline{toc}{chapter}{\textcolor{ocre}{Bibliography}}
\section*{Books}
\addcontentsline{toc}{section}{Books}
\printbibliography[heading=bibempty,type=book]
\section*{Articles}
\addcontentsline{toc}{section}{Articles}
\printbibliography[heading=bibempty,type=article]

%----------------------------------------------------------------------------------------
%	INDEX
%----------------------------------------------------------------------------------------

\cleardoublepage
\phantomsection
\setlength{\columnsep}{0.75cm}
\addcontentsline{toc}{chapter}{\textcolor{ocre}{Index}}
\printindex

%----------------------------------------------------------------------------------------

\end{document}